%% ----------------------------------------------------------------
%% Introduction.tex
%% ---------------------------------------------------------------- 
\chapter{Introduction}\label{chapter:introduction}

The recent development of low-altitude hypersonic vehicles for both civilian and military applications has renewed scientific interest in high-speed aerothermodynamics. The hypersonic regime is conventionally defined as flows above Mach~5, where strong shock-wave compressions dramatically raise the flow's stagnation temperature, creating a high-enthalpy environment that drives molecular dissociation and vibrational excitation, leading to thermochemical non-equilibrium~\citep{Candler2019}. Within the boundary layer, viscous dissipation converts kinetic energy into thermal energy, producing steep temperature gradients and elevated surface heat fluxes. From an engineering perspective, accurately predicting these effects along with associated pressure and shear stress distributions is crucial for determining aerodynamic performance and for sizing the \gls{TPS}, which must protect the structure from the resulting extreme thermal environment. \fref{figure:trajectory} compares the trajectories of various existing and experimental vehicles, including the National Aerospace Plane (NASP)~\citep{DSB1988,DSB1992}, a representative \gls{HGV}~\citep{Rizvi2015}, and the Space Shuttle Orbiter, illustrating how different flight paths generate varying thermal loads that the \gls{TPS} must be designed to withstand. The boundary layer along these trajectories is affected by real-gas effects that modify temperature and velocity profiles, as well as transport properties such as viscosity, thermal conductivity, and species diffusion. Capturing these effects accurately requires models that go beyond perfect-gas assumptions.

The heat flux on the vehicle's surface depends strongly on whether the boundary-layer flow is laminar or turbulent, with heating loads increasing by a factor of 2--10 for turbulent flow~\citep{White2005,VanDriest1952} compared to laminar flow. This adds further complexity not only for the design and sizing of the \gls{TPS} but also for predicting the aerodynamic forces and moments required to maintain stability and control throughout the entire flight envelope~\citep{Dalle2011}. The location and extent of transition to turbulence are therefore critical design considerations, yet remain highly uncertain due to the vast parameter space, the sensitivity of transition to flow conditions---including Reynolds number, Mach number, surface geometry, free-stream turbulence and acoustic noise, pressure gradients, wall roughness and porosity, boundary-layer cooling or heating, and mass injection or suction~\citep{Cebeci2001}---uncertainties in the disturbance environment, and the inherent complexity of the nonlinear processes that drive the final breakdown to turbulence. Accurate prediction of transition onset and breakdown location is essential to ensure that the \gls{TPS} is neither under- nor over-designed for mission conditions. Consequently, the use of overly conservative transition models can lead to unnecessary weight penalties in the \gls{TPS} design.

\begin{figure}[hbt!]
    \centering
    \fontsize{9}{10}\selectfont\includesvg[width=0.96\textwidth]{resources/figures/introduction/schematictrajectory.svg}
    \caption[Schematic of a small-sized waverider going at hypersonic speed in the Earth’s atmosphere experiencing real gas phenomena adapted from \citep{Anderson2019} with the HGV glide path imposed from \citep{Rizvi2015}.]{Schematic of a small-sized waverider going at hypersonic speed in the Earth’s atmosphere experiencing real gas phenomena adapted from \citep{Anderson2019} with the ~\gls{HGV} glide path imposed from \citep{Rizvi2015}.}
    \label{figure:trajectory}
\end{figure}

For a representative \gls{HGV}, downstream of the leading edge, the intense thermal and pressure gradients generated by the shock induce strong nonequilibrium effects that persist within the boundary layer. Whilst translational and rotational modes equilibrate rapidly, vibrational and chemical relaxation occur over longer timescales, resulting in a non-uniform thermochemical state along the glide path. At low altitudes and high speeds, elevated atmospheric density produces high Reynolds numbers, making the boundary layer more susceptible to transition. Even small perturbations---such as \gls{TPS} ablation, stagnation-region receptivity, or free-stream disturbances---can promote transition and substantially increase surface thermal loads~\citep{MiroMiro2019}. High stagnation enthalpies and real-gas effects further complicate this by altering boundary-layer stability characteristics~\citep{Anderson2019}. These coupled effects highlight the importance of reliable transition prediction to maintain both thermal protection and aerodynamic control without excessive design margins.


Accurately modelling the physics under such extreme conditions presents significant challenges, both experimentally and computationally~\citep{DiRenzo2021}. Only a few experimental facilities can replicate the high-enthalpy environments relevant to hypersonic flows, and even these are limited by the substantial power demands of test rigs and the need for quiet flow conditions in transition studies. As a result, most research has focused on the early stages of transition, whilst a comprehensive understanding of the breakdown pathways by which instability modes evolve into turbulence remains incomplete for hypersonic boundary layers. A deeper understanding of boundary-layer behaviour at high Mach numbers and low altitudes---where elevated temperatures, turbulence, and real-gas effects interact---is crucial for safe and efficient vehicle design. These aspects remain insufficiently explored, highlighting the need for studies that address their coupled effects.

\section{High-enthalpy flows}
In a high-enthalpy environment, gas molecules undergo complex physicochemical processes that couple macroscopic flow properties more strongly with the microscopic internal energy modes than in classical low-enthalpy flows. These internal modes---including translation, rotation, vibration, dissociation, electronic excitation, radiation, and ionisation---absorb, store, and exchange energy at quantised levels~\cite[ch.~11]{Anderson2019}. At sufficiently high temperatures, vibrational and chemical modes can lag behind the translational--rotational modes and exhibit different relaxation times, giving rise to thermochemical nonequilibrium. This delayed relaxation introduces additional modelling challenges due to the intricate coupling between thermal and chemical processes.

In hypersonic flows, internal energy modes are predominantly endothermic, acting as heat sinks that absorb energy unlike exothermic flows such as combustion, where energy is released. The accessibility of each mode depends on the enthalpy associated with the vehicle's flight conditions~\citep{Gupta1990}, with each mode requiring a threshold energy before becoming active. As temperature increases, the internal behaviour of gas species changes significantly; a detailed depiction of this behaviour in the Earth's atmosphere is illustrated in \fref{figure:tempranges}. In conventional flows (i.e. below approximately 800~K), vibrational modes remain largely inactive, and the internal energy is governed primarily by translational and rotational contributions. Under these conditions, the gas can be treated as a~\gls{CPG}, with constant specific heats and a constant ratio of specific heats, $\gamma$, more appropriately known as \textit{cold flow} in the literature. Above 800~K, vibrational excitation begins to influence thermodynamic behaviour, causing the specific heats $c_p$ and $c_v$ to vary with temperature. The gas then behaves as \emph{thermally perfect} but \emph{calorically imperfect} because $\gamma$ varies with temperature. At elevated temperatures, chemical reactions become significant, driving the dissociation and ionisation of atmospheric species over specific temperature ranges, as illustrated in \fref{figure:tempranges}~\citep{Scoggins2020}. For example, molecular oxygen (\ce{O2}) begins to dissociate at around 2500~K and is almost fully dissociated by 4000~K. These processes increase the thermochemical complexity of the flow.


\begin{figure}[h!]
    \centering
    \includegraphics[width=12cm]{resources/figures/introduction/schematic_tempranges.pdf}
    \caption[Temperature ranges at which vibrational excitation, dissociation and ionisation activate taken for air at 1-atm. Adapted from \cite{Anderson2019}.]{Temperature ranges at which vibrational excitation, dissociation and ionisation activate taken for air at 1-atm. Adapted from \cite{Anderson2019}. In general, the dissociation, and ionisation temperatures are directly proportional to the pressure.}
    \label{figure:tempranges}
\end{figure}


At high temperatures, additional molecular energy modes such as vibrational and electronic excitation (see \fref{figure:tempranges}) necessitate a multi-temperature model to capture the relevant flow physics. Gas molecules store energy in distinct translational, rotational, vibrational, and electronic modes, each with different characteristic energy levels. These are often represented by separate temperatures: $T_t$ for translation, $T_r$ for rotation, $T_v$ for vibration, and $T_e$ for electronic excitation. A multi-temperature approach models these modes individually, providing a more accurate description of the internal energy distribution. A common simplification is to assume that the translational and rotational modes behave similarly, combining them into a single translational--rotational temperature, $T_{tr}$. When all modes are considered, the vibrational and electronic modes are often assumed to equilibrate with each other~\citep{Bertin2006}, resulting in a single vibrational temperature, $T_v$. This forms the basis of Park's classic two-temperature model~\citep{Park1990}, which remains widely used for practical high-temperature flow simulations. Although multi-temperature models can match experimental shock-layer results more accurately, as shown by~\citet{Gulhan2018}, other studies suggest that detailed mode separation is only significant for expansion nozzle flows or extreme nonequilibrium conditions~\citep{ParkLee1995}. The prevailing consensus is that the increased computational cost of solving additional temperature equations is typically not justified in routine engineering practice, except in flows exhibiting pronounced nonequilibrium effects~\citep{Nagnibeda2009}. Regardless, modelling high-enthalpy effects with sufficient fidelity is essential for predicting hypersonic flow behaviour in regimes where thermal and chemical nonequilibrium strongly influence shock-layer structure, boundary-layer stability, and transition to turbulence. Capturing the correct physics and applying appropriate flow modelling strategies are therefore critical for achieving accurate and physically meaningful results.

\subsection{Flow definitions}
Chemical and vibrational processes occur over finite timescales due to molecular collisions. The relevant timescales are the vibrational relaxation timescale $\tau_{\mathrm{vib}}$, the chemical reaction timescale $\tau_{\mathrm{chem}}$, and the flow or residence timescale $\tau_{\mathrm{flow}}$. The relative magnitudes of these timescales give rise to distinct flow regimes, each with important implications for high-enthalpy flow behaviour. These relationships influence the rate of energy transfer between different modes, affect the onset and progression of chemical reactions, and shape the flow's overall thermal and chemical structure. The behaviour of the flow can be characterised using the Damköhler number, which compares the rate of a thermochemical process to the rate of convection. It is defined as the ratio of a characteristic chemical or vibrational timescale to the characteristic flow or residence timescale. For thermochemical flows, the relevant dimensionless groups are the chemical Damköhler number \gls{Dachem} and the vibrational Damköhler number \gls{Davib}, defined respectively as~\citep{miromiroThesis}:
\begin{equation}\label{equation:davib}
    \mathrm{Da}_{\mathrm{vib}} = \cfrac{\tau_{\mathrm{flow}}}{\tau_{\mathrm{vib}}}
\end{equation}
\begin{equation}\label{equation:dachem}
    \mathrm{Da}_{\mathrm{chem}} = \cfrac{\tau_{\mathrm{flow}}}{\tau_{\mathrm{chem}}}
\end{equation}
Depending on the relative magnitudes of the vibrational, chemical, and flow timescales, the flow regime may fall into several distinct categories. In the most general case, no assumptions are made about the relative ordering of the relevant timescales, giving rise to the \gls{TCNE} regime. All modes are of comparable magnitude, i.e. \gls{Davib} $= \mathcal{O}(1)$ and \gls{Dachem} $= \mathcal{O}(1)$. Considering the limiting case of \gls{TCE}, it is assumed that all internal modes---translational, rotational, vibrational, and chemical---adjust instantaneously relative to the flow timescales. This implies that the vibrational and chemical Damköhler numbers are much greater than one, i.e. \gls{Davib} $\gg 1$ and \gls{Dachem} $\gg 1$. For a \gls{TCF} case, vibrational and chemical modes do not have sufficient time to adjust during the flow process. As a result, the vibrational and chemical Damköhler numbers are much smaller than one, i.e. \gls{Davib} $\ll 1$ and \gls{Dachem} $\ll 1$. Isolating the Damköhler numbers allows for intermediate regimes to be examined. For example, if chemical reactions are frozen whilst vibrational modes adjust rapidly, then \gls{Dachem} $\ll 1$ and \gls{Davib} $\gg 1$. Conversely, if vibrational relaxation is effectively frozen whilst chemical reactions proceed, then \gls{Davib} $\ll 1$ and \gls{Dachem} $\gg 1$. These limiting cases provide insight into the individual roles of chemical and vibrational processes in shaping the overall flow behaviour.

\section{Boundary-layer transition}
\subsection{Paths to transition}
The transition from a laminar to a turbulent boundary layer is a critical aerodynamic phenomenon with significant implications for skin friction, heat transfer, and aerodynamic drag. Extensive research has sought to understand, predict, and control this process through both theoretical and experimental investigations. Within this broader context, \citet{Morkovin1994} provided a conceptual framework that synthesised decades of developments into the modern understanding of the multiple transition paths that can arise in boundary-layer flows. \fref{figure:transitionmarkovin}, reproduced from \citet{Fedorov2011}, presents a schematic overview of these paths to transition, which is applicable to both low-speed and high-speed boundary layers, although the specific instability mechanisms differ between the two regimes.

\begin{figure}[hbt!]
    \centering
    \fontsize{9}{10}\selectfont\includesvg[width=0.50\textwidth]{resources/figures/introduction/schematic_fedorov.svg}
    \caption{\cite{Morkovin1994} roadmap to transition, diagram of pathways by which a boundary layer might transition to turbulence \citep{Fedorov2011}.}
    \label{figure:transitionmarkovin}
\end{figure}

Transition begins with the appearance and amplification of disturbances originating from environmental sources such as free-stream turbulence, acoustic noise, surface roughness, or shock interactions. These external influences interact with the boundary layer and generate oscillatory modes whose amplitude, frequency, and phase define the initial conditions for instability waves that grow and eventually break down into turbulence. This conversion process, known as receptivity, governs how external disturbances give rise to eigenmode instabilities, which amplify linearly before undergoing nonlinear interactions and secondary instabilities---a sequence formalised in Morkovin's transition roadmap~\citep{Morkovin1994}.

The different pathways to turbulence are illustrated schematically in \fref{figure:transitionmarkovin}. In the most general case, low-amplitude oscillatory modes are triggered within the boundary layer and amplify through linear modal growth, followed by parametric instabilities and mode interactions that transfer energy between modes and lead to breakdown to turbulence (\textit{path A}). At higher disturbance levels, the superposition of perturbations may initially exhibit algebraic transient growth before feeding into modal amplification (\textit{path B}). As the disturbance amplitude increases further, transient growth can directly trigger parametric instabilities and mode interactions or, if sufficiently strong, bypass these stages entirely and transition straight to breakdown (\textit{paths C} and \textit{D}). Amongst these, only \textit{paths A}, \textit{B}, and \textit{C} are typically relevant for external hypersonic boundary layers, where disturbances evolve through receptivity, transient growth, and modal amplification~\citep{Lefieux2021}. In contrast, \textit{path D} is more representative of internal flows with elevated turbulence levels, whilst \textit{path E} corresponds to extreme forcing conditions---such as intense free-stream turbulence or severe surface roughness---where linear mechanisms are completely bypassed~\citep{Reshotko2001}. These physical insights have informed the development of a range of transition prediction and control methods, from empirical correlations to physics-based linear stability theory and modern numerical simulations, each aiming to capture the complex processes through which instabilities grow and interact.

In hypersonic boundary layers, key destabilising factors include free-stream perturbations (acoustic, entropy, and vorticity fluctuations), localised surface roughness, shock-wave interactions, and high-temperature effects such as vibrational excitation and dissociation, all of which can significantly enhance instability growth~\citep{vandereynde2015,BitterShepherd2015}. The path to transition is strongly influenced by both the amplitude and spectral content of these disturbances, making receptivity---the process by which external perturbations generate internal instability modes---a critical area of study. Due to the broadband and complex nature of environmental inputs, the receptivity problem remains largely unresolved in compressible boundary layers despite extensive research~\citep{Cerminara2019}. Experimental studies are particularly challenging owing to the stringent requirement for ultra-quiet wind tunnel environments~\citep{Fedorov2011}. As a result, theoretical frameworks have been widely explored, with~\citet{Fedorov2001,Fedorov2002,MaZhong2003a,MaZhong2003b} analysing receptivity mechanisms in cold hypersonic flow over a flat plate; \citet{Fedorov2002} in particular investigated receptivity to wall-induced disturbances across various wave vectors. Further studies on broadband receptivity over flared cones~\citep{Balakumar2018,HeZhong2022} showed that disturbances modified by shock interactions exhibit altered characteristics. Recently, \citet{Varma2025} examined receptivity to acoustic disturbances under thermochemical nonequilibrium and found it significantly impacts instability growth and transition onset location. Nonetheless, the concept of receptivity remains widely underexplored, and as noted by \citet{Zhong2012}, the dominant instability mechanisms differ between low- and high-speed regimes, underscoring the need for accurate physical modelling to reliably predict transition.

\subsection{Modal growth}
The transition process has been a subject of interest since the late 19\textsuperscript{th} century, beginning with the work of Lord Rayleigh~\citep{Rayleigh1878,Rayleigh1880}, who postulated the existence of sinusoidal disturbances within the incompressible boundary layer prior to transition. This mathematical framework was extended by \citet{Tollmien1936}, who showed that incompressible flows are stable to inviscid, temporally growing disturbances unless the two-dimensional mean velocity profile contains an inflection point, $\partial_y (\partial_y U) = 0$, which provides the necessary shear layer for inviscid instabilities to extract energy from the mean flow. This theoretical prediction was confirmed experimentally by \citet{Schubauer1948}, who provided clear evidence of downstream disturbance amplification using hot-wire measurements. When the effect of viscosity was considered, these unstable waves formed the classic Tollmien--Schlichting waves~\citep{Schlichting2000}. Building on this, \citet{LeesLin1946} were the first to investigate instability propagation in compressible boundary layers, although they considered only a subset of possible modes. They demonstrated theoretically that a two-dimensional, inviscid compressible boundary layer remains stable unless the mean velocity and density profiles contain a generalised inflection point, given by the condition $\partial_y ( \rho \partial_y u) = 0$, thereby extending the classic inviscid stability criterion of \citet{Tollmien1936}. They further distinguished between subsonic ($u_e - c_r < a_e$), sonic ($u_e - c_r = a_e$), and supersonic ($u_e - c_r > a_e$) disturbances, where $c_r$ is the phase speed of the disturbance, and $u_e$ and $a_e$ are the streamwise velocity and speed of sound at the boundary-layer edge, respectively. This theoretical prediction was later supported by experimental observations from \citet{Laufer1960}, who helped constrain the stability boundaries for compressible flows. They showed that an adiabatic compressible boundary layer is more stable than its incompressible counterpart.

A rigorous mathematical framework was first developed by \citet{LeesReshotko1962} to analyse two-dimensional subsonic disturbances, marking a significant departure from earlier asymptotic approaches. Their analysis showed that, although the transition process retains structural similarities with incompressible flows, the emergence of multiple instability loops and the increasing influence of viscous dissipation at high Mach numbers point to fundamentally different dominant mechanisms in supersonic regimes. A comprehensive \gls{LST} analysis was later performed by \citet{Mack1963,Mack1964}, who clearly distinguished between different instability modes. The compressible analogue of Tollmien--Schlichting waves, initially observed by \citet{LeesLin1946}, was identified as the first mode, whilst additional high-frequency acoustic disturbances emerging at higher Mach numbers were classified as distinct instabilities. Due to the hyperbolic nature of the compressible governing equations in supersonic and hypersonic flows, an infinite set of eigenvalues arises in the stability problem; the most unstable of these, identified by Mack, became known as the second mode~\citep{Hudson1997}. \citet{Mack1975} further demonstrated that the first mode becomes increasingly stabilised with rising Mach number and decreasing wall temperature, thereby rendering the second mode the dominant instability mechanism in high-speed boundary layers.


At low Mach numbers, where the flow is effectively incompressible, transition is primarily governed by two-dimensional Tollmien--Schlichting waves, which are vorticity-dominated, destabilised by viscous effects, and prone to amplification or modal interactions in three-dimensional configurations~\citep{Schubauer1948}. The first mode emerges in the transonic and early supersonic regime, typically as a three-dimensional oblique instability driven by cross-stream perturbation coupling and reduced upstream influence~\citep{Mack1984,Reed1996}. At higher Mach numbers, the second mode, commonly referred to as the \textit{Mack mode} becomes dominant~\citep{Kuehl2018}. This instability arises from acoustic waves trapped between the wall and the relative sonic line above the boundary layer~\citep{Bitter2015}, with disturbances convected subsonically above the sonic line and supersonically beneath it. In two-dimensional boundary layers, the most amplified second-mode disturbances remain strictly two-dimensional, whereas in three-dimensional flows, the most unstable waves exhibit slight obliqueness~\citep{Balakumar1991}. However, \citet{Franko2013} observed that three-dimensional Mack modes undergo delayed transition relative to their two-dimensional counterparts in zero-pressure-gradient boundary layers, suggesting that the mechanisms governing their amplification remain incompletely understood. A distinct class of additional modes, referred to as \textit{supersonic modes}, was later identified by \citet{Mack1987}. These share some characteristics with the second mode but propagate supersonically with respect to the edge velocity. Supersonic modes are commonly observed in flows over highly cooled walls and, unlike the previously discussed instabilities, radiate energy into the free stream, resulting in a perturbation amplitude that does not decay exponentially with wall-normal distance. Further modes have also been identified in configurations with pressure gradients, where crossflow instabilities and Görtler vortices become prominent features of the transition process. The mechanisms governing this shift in dominant instability behaviour remain an active area of research. In particular, the interaction between crossflow modes and high-frequency acoustic disturbances has been investigated as a means of distinguishing competing transition scenarios~\citep{Dong2020}. Considering a general \gls{ZPG} flat-plate boundary layer, the prevailing consensus in the literature is that the dominant instability mode is \textit{Mack's second mode}, typically associated with frequencies in the hundreds of kilohertz range. Amongst the various instability mechanisms, these high-frequency acoustic modes are particularly relevant in hypersonic regimes due to their rapid amplification, strong wall-normal gradients, and central role in triggering transition.


% Since their inception, \textit{Mack's second modes} have been the subject of extensive numerical and experimental investigations aimed at characterising the behaviour of these high-frequency acoustic instabilities. \citet{Kendall1975} was the first to confirm the existence of the second mode experimentally. Subsequently, \citet{Lysenko1984} demonstrated that the second mode remains the dominant instability mechanism for Mach numbers above 4, whilst also highlighting the influence of wall cooling on both first and second modes. Building on the \gls{LST} framework established by \citet{Mack1963}, \citet{Erlebacher1990} conducted temporal \gls{DNS} to examine the stability characteristics of the second mode and their subsequent breakdown to turbulence, and confirmed the presence of additional acoustic instabilities. Despite decades of research, a comprehensive understanding of Mack modes, particularly their nonlinear interactions and behaviour under thermally excited, real-gas conditions remains elusive~\citep{MiroMiro2019}. The linear and early nonlinear stages, along with the role of thermochemical nonequilibrium in typical flight scenarios, therefore remain key areas of interest and will be explored in detail in later sections.

% \subsection{Breakdown to turbulence}

Since their inception, \textit{Mack's second modes} have been the subject of extensive numerical and experimental investigations aimed at characterising the behaviour of these high-frequency acoustic instabilities. \citet{Kendall1975} was the first to confirm the existence of the second mode experimentally. Subsequently, \citet{Lysenko1984} demonstrated that the second mode remains the dominant instability mechanism for Mach numbers above 4, whilst also highlighting the influence of wall cooling on both first and second modes. Building on the \gls{LST} framework established by \citet{Mack1963}, \citet{Erlebacher1990} conducted temporal \gls{DNS} to examine the stability characteristics of the second mode and their subsequent breakdown to turbulence, confirming the presence of additional acoustic instabilities. Much of this work has focused on relatively benign thermochemical conditions—calorically perfect gas or modest temperatures where real-gas effects remain limited. However, realistic flight conditions at low altitudes and high speeds involve significant enthalpy levels where vibrational excitation, chemical reactions, and thermal nonequilibrium become important. The behaviour of second-mode instabilities under these conditions, particularly the nonlinear interactions and breakdown mechanisms in thermally excited environments, remains incompletely understood~\citep{MiroMiro2019}. Understanding how modal disturbances ultimately transition to fully turbulent flow under such conditions is therefore essential, motivating detailed examination of the breakdown process.

\subsection{Breakdown to Turbulence}

Following the initial amplification of instability modes predicted by linear theory, the disturbances eventually grow to sufficient amplitude to undergo nonlinear saturation, which leads to secondary instabilities and the eventual breakdown to turbulence. Only the initial amplification phase is captured by linear theory, with the subsequent nonlinear processes requiring alternative approaches. The transition from primary to secondary instability occurs when the amplitude of the most unstable primary mode reaches approximately 1\% of the free-stream velocity~\citep{Klebanoff1962}. At this threshold, nonlinear saturation modifies the base flow to a new quasi-steady state upon which secondary disturbances can develop. The framework for analysing this evolution in incompressible flows was established by \citet{Herbert1988}, who developed secondary instability theory based on Floquet analysis. In this approach, three-dimensional disturbances of prescribed wavelength and frequency are superimposed onto the saturated primary wave, and their growth rates are determined through linear stability analysis of the modified base flow. These nonlinear interactions excite additional spatial and temporal scales, inducing harmonics and subharmonics of the dominant disturbance frequencies, thereby accelerating the breakdown of the laminar flow into the chaotic motion and broadband spectrum characteristic of turbulence. 

Experimental investigations by \citet{KachanovLevchenko1984} and \citet{SaricThomas1984} subsequently identified three distinct classes of linear secondary instabilities, alongside the characteristic lambda-shaped vortical structures observed near the boundary layer edge. The first class, designated \textit{K-type} instability after \citet{Klebanoff1962}, represents fundamental resonance wherein secondary waves share the primary wave's streamwise wavenumber whilst exhibiting spanwise periodicity aligned with the primary structure. The second class, subharmonic or \textit{H-type} instability named after \citet{Herbert1988}, features secondary disturbances with half the streamwise wavelength of the primary wave, producing a staggered vortical arrangement that has been shown to be the most dominant secondary instability mechanism in incompressible flows. The third class, \textit{O-type} instability, denotes oblique breakdown~\citep{Fasel1993}, which originates from the nonlinear interaction of two oblique instability waves with equal wave angles but in opposite directions. The aforementioned investigations established the secondary instability framework for incompressible flows, identifying a notable discrepancy wherein \textit{H-type} instabilities feature prominently in theoretical predictions but \textit{K-type} modes dominate experimental measurements~\citep{SpalartYang1987}.

% When secondary instability theory is applied to compressible regimes, significant departures from incompressible behaviour emerge. ElHady1992
Early numerical investigations of compressible boundary-layer transition employed spatial simulations with various combinations of initial disturbances \citep{ThummWolzFasel1989}. At Mach~4.5, \citet{Erlebacher1990} revealed that the coupling between two-dimensional and oblique disturbances initiates secondary instabilities that share features with the \textit{K-type} breakdown mechanism observed in incompressible flows. Crucially, their simulations confirmed that compressibility introduces a stabilising effect relative to incompressible boundary layers. Building upon this foundation, \citet{NgErlebacher1992} conducted a systematic parametric investigation spanning multiple Mach numbers, demonstrating that secondary instabilities linked to the primary Mack mode exhibit enhanced growth compared to those originating from primary Tollmien--Schlichting waves as Mach number increases. Their analysis further identified subharmonic and combination resonance modes as the predominant instability mechanisms. However, a competing trend was also observed: higher Mach numbers suppress the growth rates of subharmonic disturbances~\citep{MasadNayfeh1990,NgErlebacher1992,ElHady1992}, indicating a complex interplay between compressibility effects and secondary instability characteristics.


Through experiments and temporal-\gls{DNS}, several studies on low supersonic flows~\citep{KosinovMaslovShevelkov1990,Fasel1993,SandhamAdams1993,ChangMalik1994} found that sufficiently strong nonlinear interactions of oblique waves can bypass secondary instability stages, precipitating direct laminar--turbulent transition. \citet{SandhamAdamsKleiser1995} elucidated the mechanism: nonlinear interactions of primary three-dimensional disturbances generate streamwise vortices and high-shear layers that roll up into additional vortices, culminating in breakdown. \citet{SchmidHenningson2001} examined oblique breakdown, subharmonic transition, and fundamental transition across incompressible and supersonic regimes, establishing a comparative framework for compressibility effects.

The transition process initiates with $\Lambda$-shaped vortices near the critical layer, generated from oblique subharmonic instabilities. Flow visualisations by \citet{AdamsKleiser1996} revealed that these $\Lambda$-vortices produce \textit{Y-shaped} shear layers below and between neighbouring vortices. These layers break up to create new vortices above and below the critical layer, initiating breakdown. A secondary shear layer forms at the symmetry plane near the boundary-layer edge, whose break-up drives the final transition stages. Two distinct breakdown scenarios were documented: one where Y-shaped shear layers break up above the sonic layer absent in low Mach flows and another above the boundary-layer edge resembling incompressible transition. In agreement with \citet{NgErlebacher1992}, subharmonic secondary instability emerged as the dominant mechanism, with no evidence of fundamental-type transition.


% Through a series of experiments and temporal-\gls{DNS}, several studies on low supersonic compressible flows~\citep{KosinovMaslovShevelkov1990,Fasel1993,SandhamAdams1993,ChangMalik1994} found that primary instabilities can directly precipitate laminar--turbulent transition. When nonlinear interactions of oblique waves are sufficiently strong, they can bypass secondary instability stages and break down directly into turbulence. \citet{SandhamAdamsKleiser1995} elucidated the underlying mechanism: nonlinear interactions of primary three-dimensional disturbances generate streamwise vortices with high-shear layers. These layers subsequently roll up to produce additional vortices, eventually leading to breakdown. The sustained growth of quasi-streamwise vortices is maintained by the nonlinear interaction of oblique primary instability waves, which generate high-shear layers that roll up into additional vortices, culminating in turbulent breakdown. \citet{SchmidHenningson2001} examined oblique breakdown, subharmonic transition, and fundamental transition across both incompressible and supersonic regimes, establishing a comparative framework for understanding compressibility effects.

% The transition process initiates with the formation of $\Lambda$-shaped vortices near the critical layer, generated from oblique subharmonic instabilities. Two distinct breakdown scenarios have been documented. In the first, Y-shaped shear layers break up to create new vortices located above the sonic layer---a pattern not observed in low Mach number flows. The second scenario occurs above the boundary-layer edge and resembles the incompressible case, where $\Lambda$-vortices peak primarily at peak planes. Flow visualisations by \citet{AdamsKleiser1996} revealed that $\Lambda$-vortices generate \textit{Y-shaped} shear layers located below and between neighbouring vortices. These shear layers break up to produce additional vortices above and below the critical layer, initiating the first stages of breakdown. A secondary type of shear layer forms at the symmetry plane of the $\Lambda$-vortices near the boundary-layer edge, whose break-up drives the final stages of transition. In agreement with \citet{NgErlebacher1992}, no evidence of fundamental-type transition was observed in these compressible flows, with subharmonic secondary instability emerging as the dominant mechanism.

% The extension of secondary instability theory to compressible flows has revealed important differences from incompressible counterparts. \citet{Erlebacher1990} investigated a Mach~4.5 boundary layer and identified that interactions between two-dimensional and oblique waves trigger secondary instabilities exhibiting characteristics similar to \textit{K-type} breakdown.  They further demonstrated that compressible boundary layers are inherently more stable than their incompressible counterparts. This work was extended by \citet{NgErlebacher1992} who performed a parametric study across a range of Mach numbers, revealing that secondary instabilities associated with the primary Mack mode are stronger than those of the primary Tollmien--Schlichting wave at higher Mach numbers. Moreover, the dominant instabilities are associated with the subharmonic and the combination resonance modes.  Overall, increasing the Mach number minimises the subharmonic growth rate \citep{MasadNayfeh1990,NgErlebacher1992}.




% The extension of secondary instability theory to compressible flows has revealed important differences from incompressible counterparts. \citet{SchmidHenningson2001} examined oblique breakdown, subharmonic transition, and fundamental transition across both incompressible and supersonic regimes, establishing a comparative framework for understanding compressibility effects. \citet{Erlebacher1990} investigated a Mach~4.5 boundary layer and identified that interactions between two-dimensional and oblique waves trigger secondary instabilities exhibiting characteristics similar to \textit{K-type} breakdown. They further demonstrated that compressible boundary layers are inherently more stable than their incompressible counterparts. This work was extended by \citet{NgErlebacher1992} across a range of Mach numbers, revealing that secondary instabilities associated with the primary Mack mode are stronger than those of the primary Tollmien--Schlichting wave at elevated Mach numbers.

% At lower supersonic speeds, several studies~\citep{ChangMalik1994,Fasel1993,SandhamAdams1993} found that primary instabilities can directly precipitate laminar--turbulent transition. When nonlinear interactions of oblique waves are sufficiently strong, they can bypass secondary instability stages and break down directly into turbulence. \citet{SandhamAdamsKleiser1995} elucidated the underlying mechanism: nonlinear interactions of primary three-dimensional disturbances generate streamwise vortices with high-shear layers. These layers subsequently roll up to produce additional vortices, eventually leading to breakdown. The sustained growth of quasi-streamwise vortices is maintained by the nonlinear interaction of oblique primary instability waves, which generate high-shear layers that roll up into additional vortices, culminating in turbulent breakdown.



% Looking into the compressibility implications, \cite{SchmidHenningson2001} studied oblique breakdown, subharmonic transition and fundamental transition in both incompressible and supersonic regime. \cite{Erlebacher1990} studied compressible Mach~4.5 boundary layer and found that an interaction between a two-dimensional and an oblique wave triggers secondary instability, which showed similar traits to the \textit{K-type} breakdown scenario. They also found that compressible boundary layers are inhereintly more stable than incompressible counterparts. \cite{NgErlebacher1992} extended the studies to a range of Mach numbers, and found that secondary instability is stronger for the primary Mach mode than the primary TS-wave at higher Mach numbers. In the compressible boundary layer, at lower supersonic speeds, \citet{ChangMalik1994},\citep{Fasel1993}, \cite{SandhamAdams1993}) found that primary instabilities can be responsible for the laminar-turbulent transition. If nonlinear interactions of oblique waves are strong enough, they can directly break down into turbulence. In this case, no secondary instability is necessary. In the subsequent year, El-Hady14 and Masad and Nayfeh investigated the effect of suction and wall cooling on the secondary instability in the compressible boundarylayer.\cite{SandhamAdamsKleiser1995} ffound that the reason for
% this behaviour was nonlinear interactions of primary three-dimensional disturbances generate vortices in the streamwise direction with high shear layer. This layer will roll up and produce more vortices, eventually breaking down into turbulence. It was found that the nonlinear interaction of oblique primary instability waves sustains the growth of quasi-streamwise vortices. These vortices generate high shear layers which rollup into additional vortices leading to breakdown to turbulence. The transition process starts from the formation of
% lambda-vortices near the critical layer generated from oblique subharmonic instabilities. Two different transition scenarios were documented. In the first case, a Y-shaped shear layer breaks up to create new vortices located above the sonic layer. The second phase is located above the boundary layer edge and is similar to the incompressible case, where the lambda-vortices peaks mainly at peak planes. This transition pattern is not observed in the low Mach number flows. In agreement with \citet{NgErlebacher1992}, no evidence of the fundamental type of transition was found, so attention was focused on the subharmonic secondary instability. Flow visualisations revealed the formation of $\Lambda$-shaped vortices close to the boundary-layer edge. The $\Lambda$-vortices generate Y-shaped shear layers, located below and between two neighbouring vortices~\citep{AdamsKleiser1996}, which break up to produce additional vortices above and below the critical layer, thereby being responsible for the first stages of breakdown. A second type of shear layer is created by the $\Lambda$-vortices at their symmetry plane and close to the boundary-layer edge. The break-up of these shear layers is responsible for the last stages of breakdown.

% Alternative pathways to breakdown exist in flows where transient growth mechanisms dominate. The evolution and breakdown of localised three-dimensional disturbances in Poiseuille flow were mapped in detail by \citet{Henningson1993}, demonstrating that small disturbances undergo substantial transient growth through the \textit{lift-up} mechanism~\citep{Landahl1975}, form elongated streamwise streaks, and subsequently evolve nonlinearly towards breakdown. The streak-induced, spanwise-modulated shear layers that emerge at finite amplitude support secondary instabilities of the streaks themselves, proliferating scales and precipitating transition through bypass mechanisms~\citep{ButlerFarrell1991}. \citet{Hanifi1996} extended both the secondary instability and transient growth frameworks to compressible flows through temporal stability analyses, establishing the foundation for understanding breakdown mechanisms in high-speed boundary layers.



% Alternative pathways to breakdown exist in flows where transient growth mechanisms dominate. \citet{Henningson1993} demonstrated that localised three-dimensional disturbances in Poiseuille flow undergo substantial transient growth through the \textit{lift-up} mechanism~\citep{Landahl1975}, forming elongated streamwise streaks that evolve nonlinearly towards breakdown through bypass transition~\citep{ButlerFarrell1991}. \citet{Hanifi1996} extended these frameworks to compressible flows. However, the present work focuses on modal secondary instabilities, which represent the dominant breakdown mechanism in low-disturbance environments.


% Following the initial amplification of instability modes predicted by linear theory, the disturbances eventually grow to sufficient amplitude to undergo nonlinear saturation, which leads to secondary instabilities and the eventual breakdown to turbulence. Only the initial amplification phase is captured by linear theory, with the subsequent nonlinear processes requiring alternative approaches. 

% The transition from primary to secondary instability occurs when the amplitude of the most unstable primary mode reaches approximately 1\% of the free-stream velocity~\citep{Klebanoff1962}. At this threshold, nonlinear saturation modifies the base flow to a new quasi-steady state upon which secondary disturbances can develop. The framework for analysing this evolution in incompressible flows was established by \citet{Herbert1988}, who developed secondary instability theory based on Floquet analysis. In this approach, three-dimensional disturbances of prescribed wavelength and frequency are superimposed onto the saturated primary wave, and their growth rates are determined through linear stability analysis of the modified base flow. These nonlinear interactions excite additional spatial and temporal scales, inducing harmonics and subharmonics of the dominant disturbance frequencies, thereby accelerating the breakdown of the laminar flow into the chaotic motion and broadband spectrum characteristic of turbulence.

% Experimental investigations by \citet{KachanovLevchenko1984} and \citet{SaricThomas1984} subsequently identified three distinct classes of linear secondary instabilities, alongside the characteristic lambda shaped vortical structures observed near the boundary layer edge. The first class, designated \textit{K-type} instability after \citet{Klebanoff1962}, represents fundamental resonance wherein secondary waves share the primary wave's streamwise wavenumber whilst exhibiting spanwise periodicity aligned with the primary structure. The second class, subharmonic or \textit{H-type} instability named after \citet{Herbert1988}, features secondary disturbances with half the streamwise wavelength of the primary wave, producing a staggered vortical arrangement that has been shown to be the most dominant secondary instability mechanism in incompressible flows. The third class, \textit{O-type} instability, denotes oblique breakdown~\citep{Fasel1993}, which originates from the nonlinear interaction of two oblique instability waves with equal wave angles but in opposite directions.

% Alternative pathways to breakdown exist in flows where transient growth mechanisms dominate. The evolution and breakdown of localised three-dimensional disturbances in Poiseuille flow were mapped in detail by \citet{Henningson1993}, demonstrating that small disturbances undergo substantial transient growth through the \textit{lift-up} mechanism~\citep{Landahl1975}, form elongated streamwise streaks, and subsequently evolve nonlinearly towards breakdown. The streak-induced, spanwise-modulated shear layers that emerge at finite amplitude support secondary instabilities of the streaks themselves, proliferating scales and precipitating transition through bypass mechanisms~\citep{ButlerFarrell1991}. \citet{Hanifi1996} extended both the secondary instability and transient growth frameworks to compressible flows through temporal stability analyses, establishing the foundation for understanding breakdown mechanisms in high-speed boundary layers.


% Following the initial amplification of instability modes predicted by linear theory, the disturbances eventually grow to sufficient amplitude to undergo nonlinear saturation~\citep{ButlerFarrell1991}, which triggers secondary instabilities and the eventual breakdown to turbulence. These nonlinear interactions excite additional spatial and temporal scales, inducing harmonics and subharmonics of the dominant disturbance frequencies, thereby accelerating the breakdown of the laminar base flow into the chaotic motion and broadband spectrum characteristic of turbulence. The evolution and breakdown of localised three-dimensional disturbances in Poiseuille flow were first mapped in detail by \citet{Henningson1993}, demonstrating that small disturbances undergo substantial transient growth through the \textit{lift-up} mechanism~\citep{Landahl1975}, form elongated streamwise streaks, and subsequently evolve nonlinearly towards breakdown. The streak-induced, spanwise-modulated shear layers that emerge at finite amplitude support secondary instabilities, which couple the primary content with harmonics and subharmonics, proliferate scales, and precipitate transition. \cite{Hanifi1996} extended the previous frameworks to compressible regime, looking into temporal transient growth analyses. 

% The transition from primary to secondary instability occurs when the amplitude of the most unstable primary mode reaches approximately 1\% of the free-stream velocity~\citep{Klebanoff1962}. At this threshold, nonlinear saturation modifies the base flow to a new steady state upon which secondary disturbances can develop. The framework for analysing this evolution in incompressible flows was established by \citet{Herbert1988}, who developed secondary instability theory based on Floquet analysis. In this approach, three-dimensional disturbances of prescribed wavelength and frequency are superimposed onto the saturated primary wave, and their growth rates are determined through linear stability analysis of the modified base flow. Experimental investigations by \citet{KachanovLevchenko1984} and \citet{SaricThomas1984} subsequently identified three distinct classes of linear secondary instabilities, alongside the characteristic lambda-shaped vortical structures observed near the boundary-layer edge. 

% The first class, designated \textit{K-type} instability after \citet{Klebanoff1962}, represents fundamental resonance. This mode is characterised by secondary waves sharing the primary wave's streamwise wavenumber whilst exhibiting spanwise periodicity aligned with the primary structure. The second class, subharmonic or \textit{H-type} instability named after \citet{Herbert1988}, features secondary disturbances with half the streamwise wavelength of the primary wave, producing a staggered vortical arrangement that has been shown to be the most dominant secondary instability mechanism in incompressible flows. The third class, \textit{O-type} instability, denotes oblique breakdown~\citep{Fasel1993}, which originates from the nonlinear interaction of two oblique instability waves with equal wave angles but in opposite directions. 

% The compressible 





% Following the initial amplification of instability modes, the disturbances grow to sufficient amplitude to undergo nonlinear saturation~\citep{ButlerFarrell1991}, which leads to secondary instabilities and the eventual breakdown to turbulence. Only the initial amplification phase is captured by linear theory, with the subsequent nonlinear processes requiring alternative approaches. These interactions excite additional spatial and temporal scales, inducing harmonics and subharmonics of the dominant disturbance frequencies, thereby accelerating breakdown of the laminar base flow into the chaotic motion and broadband spectrum characteristic of turbulence. The evolution and breakdown of localised three-dimensional disturbances in Poiseuille flow were first mapped in detail by \citet{Henningson1993}, demonstrating that small disturbances undergo substantial transient growth through the lift-up mechanism~\citep{Landahl1975}, form elongated streamwise streaks, and subsequently evolve nonlinearly towards breakdown. The streak-induced, spanwise-modulated shear layers that emerge at finite amplitude support secondary instabilities, which couple the primary content with harmonics and subharmonics, proliferate scales, and precipitate transition. When the amplitude of the most unstable primary mode reaches approximately 1\% of the free-stream velocity~\citep{Klebanoff1962}, nonlinear saturation occurs and the base flow is modified to a new steady state. The evolution from primary to secondary instability in incompressible flows was analysed by \citet{Herbert1988}, who developed the secondary instability framework. In this approach, three-dimensional disturbances of prescribed wavelength and frequency are superimposed onto the saturated primary wave, and Floquet analysis is used to determine the growth rate of the secondary mode. \citet{KachanovLevchenko1984} and \citet{SaricThomas1984} experimentally identified three distinct classes of linear secondary instabilities, alongside the characteristic lambda-shaped vortical structures observed near the boundary-layer edge. 

% Fundamental resonance, designated \textit{K-type} instability after \citet{Klebanoff1962}, represents the first class. This mode is characterised by secondary waves sharing the primary wave's streamwise wavenumber whilst exhibiting spanwise periodicity aligned with the primary structure. The second class, subharmonic or \textit{H-type} instability named after \citet{Herbert1988}, is characterised by secondary disturbances with half the streamwise wavelength of the primary wave, producing a staggered vortical arrangement. These were shown to be the most dominant secondary instability in incompressible flows. The third class is the \textit{O-type} instability, where \textit{O} denotes oblique breakdown~\citep{Fasel1993}. This breakdown originates from the nonlinear interaction of two oblique instability waves with equal wave angles but in opposite directions.







% The study of transition breakdown in hypersonic flows has evolved significantly since the early investigations of high-speed boundary layers. Early experimental work by \cite{Demetriades1974, Kendall1975} identified the presence of second-mode instabilities and their role in hypersonic transition. The concept of secondary instability as a precursor to breakdown was established through the pioneering theoretical work of \cite{Herbert1988} for incompressible flows and later extended to compressible regimes by \cite{Ng1999, Mayer2011}. Direct numerical simulations (DNS) have provided unprecedented insight into the breakdown process, with landmark studies by \cite{Pruett1995, Sandham2014} revealing the complex three-dimensional structures that emerge during the final stages of transition.


% Following the initial amplification of instability modes predicted by linear theories---and, where relevant, transient non-modal growth~\citep{ButlerFarrell1991,Trefethen1993,SchmidHenningson2001}---their amplitudes eventually become sufficient to distort the laminar base flow and trigger nonlinear interactions amongst modes. These interactions excite additional spatial and temporal scales, inducing harmonics and subharmonics of the dominant disturbance frequencies, thereby accelerating breakdown of the laminar base flow into the chaotic motion and broadband spectrum characteristic of turbulence. The evolution and breakdown of localised three-dimensional disturbances in Poiseuille flow were first mapped in detail by \citet{Henningson1993}, demonstrating that small disturbances undergo substantial transient growth through the lift-up mechanism~\citep{Landahl1975}, form elongated streamwise streaks, and subsequently evolve nonlinearly towards breakdown. The streak-induced, spanwise-modulated shear layers that emerge at finite amplitude support secondary instabilities, which couple the primary content with harmonics and subharmonics, proliferate scales, and precipitate transition.

% ~\citep{Herbert1988,Kachanov1994,ReedSaricArnal1996}.


% Following the initial amplification of instbaility modes predicted by linear theories 


% Following the initial amplification of instability modes predicted by linear theories—and, where relevant, transient non-modal growth \citep{ButlerFarrell1991,Trefethen1993,SchmidHenningson2001} their amplitudes eventually become sufficient to distort the laminar base flow and trigger nonlinear interactions among modes. These interactions excite additional spatial and temporal scales, inducing harmonics and subharmonics of the dominant disturbance frequencies, thereby accelerating breakdown of the laminar base flow into the chaotic motion and broadband spectrum characteristic of turbulence. The evolution and breakdown of localised three-dimensional disturbances in Poiseuille flow were first mapped in detail by \citet{Henningson1993}, showing that small disturbances undergo substantial transient growth by the lift-up mechanism \citep{Landahl1975}, form elongated streamwise streaks, and then evolve nonlinearly to breakdown. The streak-induced, spanwise-modulated shear layers that arise at finite amplitude support secondary instabilities, which couple the primary content with harmonics and subharmonics, proliferate scales, and precipitate transition \citep{Herbert1988,Kachanov1994,ReedSaricArnal1996}. 


% Following the initial amplification of instability modes predicted by linear theories, relevant transient non-modal growth~\citep{ButlerFarrell1991} their amplitudes eventually become sufficient to distort the laminar base flow and trigger nonlinear interactions among modes. These interactions excite adiditional spatial and temporal scales, including harmonics and subharmonics of the dominant disturbance frequencies, thereby accelerating breakdown 

% References:
% \citep{ButlerFarrell1991}, 
% \citep{Trefethen1993}
% \citep{SchmidHenningson2001}
% \citep{Hanifi1996}
% \citep{Andersson2001}
% \citep{BrandtHenningson2002}
% \citep{AlamSandham2000}



% \cite{Hanifi1996} extended the previous frameworks to compressible regime, looking into temporal transient growth analyses. 

% . The finite amplitude localised disturbances were shown to undergo transition through the lift-up phenomenon \citep{Landahl1975}, 


% The evolution of a three-dimensional localized disturbance in channel flow was first examined by Henningson, Lundbladh \& Johansson (1993), who showed that a small-amplitude disturbance undergoes linear transient growth via the lift-up mechanism \citep{Landahl1975}, whereby the disturbance three-dimensionality generates energetic streamwise streaks, amplifies wall-normal perturbation vorticity, and ultimately preconditions the flow for secondary instability and breakdown.




% As the disturbances amplify according to the above mentioned mechanisms, their amplitude eventually becomes large enough to modufy the laminar flow profile and cause non-linear interactions between modes. Because of these nonlinear interactions, additional legnth and time scales are excited at harmonics of the dominant disturbance frequencies and the flow rapidly breaks down into the chaotic motion and broad-band frequency spectrum of turbulence. Path E in \fref{figure:transitionmarkovin} accounts for disturbances that are initially of such large amplitude that no amplification is needed in order to reach the onlinear regime, which is often called bypass transition. Although initial efforts to predict transition to turbulence focused on the modal instability mechanisms, it was soon realised that modal instabilities could not explain all the phenomenon observed in experiments. In particular, certain canonical flows, such as pipe flow and Couette flow, have no unstable modes, and yet transition to turbulence is readily achieved in experiments. Other flows, like plane Poiseuille flow, transition at low Reynolds numbers for which the linear theory predicts that the flow is stable. 

% The modal analysis 

% \cite{Herbert1988} talks about incompressible secondary instabilities and breakdown scenarios that include subharmionic and fundamental secondary instabilities. 

% \cite{SchmidHenningson2001} talks about transition in three different transition mechanisms for incompressible and supersonic flow. Oblique breakdown starts with the interaction of two oblique waves (either first mode or second mode), which interacts non-linearly to produce a dominant spanwise varying stationary mode. This mode contains large streamwise vorticity and leads to velocity and temperature streaks in the boundary layer. Higher modes are also generated nonliearly and the boundary layer completely transitions to turbulence. Subharmonic instability begins with the growth of a 2D wave, eitehr a first mode or a second mode instability. Onbce the 2D wave is large enough, an obliquye subharmonic mode ( frequency half of primary mode and with a spanwise wavenumber) grows rapidly. The subharmonic modes then ineract to produce a stationary mode ( similar to oblique transition). Fundamental transition is similar to
% subharmonic transition, except a three dimensional disturbance with the same frequency as the primary disturbance is generated. 


% In a boundary layer flow, after the initial transient stage is covered by modal growth. The precursor of transition to turbulence is given by the secondary instability or nonmodal growth that breaakdowns to turbulence.  The transition process over a flat plate boundary alyer is shown in the figure below. This scenario was investigated almost simultaneously in the theoretical work of Herbert  the experimetnal work of Kachavov and Levchenko and DNS works of Kleiser Schumann, Henningson. When the most unstable of primary instability grows in amplitude, for an incomrpessible flow, this is tollmein schiting waves need to reach 1\% of the  frestream. This new baseflow can be unstable to exponentially small disturbances, particularly three-dimensional waves in the boundary layer Herbert 1988. Since the new formulation will involved equations with periodic coefficients, Floquet theory analyses this newly modified state. Physically, secondary instabilities represent lambda-shaped voritces formed during the late stages of transiiton near the boundary layer edge. In this matter, three clasees of linear secondary instabilities have been idnetified experimentally by Kachanov and Levchenko and Saric and Thomas. 


% The first fundamental instability/resonance, aslo referred to as K-type instbailiity due to the experimental work of Klebanoff. In this case, the secondary waves have the exact streamwise wavenumber as the primary waves and are aligned with each other. The second one is subharmonic resonance also known as H-type instabilitye ( Herbert 1988). In this case, the streamwise wavenumber of the secondary instability is twice the primary and forms a staggered vortices pattern. DNS simulations performed by Spalart and Yang 1987 showed that H type instbaility is the most unstable secondary instability in an incompressible boundary alyer. However, K-type has been found to be dominant in experiments. This explained due to low amplitudes streamwise voriticity in the background flow. The third class is referred to as O-type instbaility (oblique breakdown) and was discovered by Thummm, Fasel. This breakdown starts from the nonlinear interaction of two oblique instability waves with equal wave angles but in opposite directions. 

% Now, let's chat about the compressible trnasition mechanism, there will be the extra things with Y-shaped transition process and it was studies in the past by Sandham and Adamas, Sandham et al, Adamns and Kleiser, Pruett and Zhang, Mayer - oblique breakdown. For maintly cold  wall flows, recently studies of 

% Due to the difficulties of experimental studies of hypersonic boundary layer transition, DNS can give insight into the potential mechanisms for breakdown to turbulence and the overshoot in wall heat transfer. This insight can then be used to guide further exploration in both computations and experiments

% \newpage

% Strictly speaking hypersonic vehicles are divided into four types~\citep{Hirschel2015}, winged reentry vehicles, hypersonic cruise vehicles, ascent and reentry vehicles and aero-assisted orbit transfer vehicles. Each type of vehicle has its own design challenges and requirements. 

% \section{Recent Progress in Hypersonics}\label{section:introduction/recentprogress}


% The hypersonic research landscape has experienced a resurgence driven by strategic competition and technological breakthroughs in propulsion, materials, testing capabilities, and computational simulation. Advances in air-breathing propulsion, ultra-high temperature ceramics, ground-based testing infrastructure, and computational fluid dynamics have fundamentally altered the feasibility of sustained atmospheric flight exceeding Mach~5. However, this progress has revealed the profound complexity of hypersonic flow physics, particularly boundary layer transition, thermochemical nonequilibrium effects, and material oxidation at extreme temperatures. Central to understanding these phenomena is boundary layer transition, which determines whether heating rates remain four to seven times lower in laminar versus turbulent flow—a distinction that fundamentally drives \gls{TPS} requirements and vehicle performance. Accurate prediction and experimental characterisation of transition under realistic hypersonic conditions therefore represents a critical research priority. Despite advances in ground-based facilities, high-enthalpy conditions representative of low-altitude, high Reynolds number flows require excessive power consumption, making simultaneous replication of all flight parameters infeasible~\citep{Bertin2006}. Nevertheless, several facilities worldwide exist to study both low- and high-enthalpy hypersonic flows across different flight regimes~\citep{Grossir2015}.


% \section{Recent Advances in Hypersonic Flow Simulation}\label{section:introduction/recent_advances}



% Recent progress in understanding modal instability growth and breakdown to turbulence in hypersonic boundary layers has come from experimental measurements and numerical simulations. Facility-based studies have documented second-mode wave development under various flow conditions, whilst computational work has examined instability mechanisms at resolutions beyond experimental capability. This research has identified key phenomena including nonlinear wave interactions, thermochemical influences on growth rates, and pathways to turbulent flow. Questions remain regarding transition behaviour where high enthalpy, finite-rate chemistry, and geometry effects combine.


% The \gls{NASA} Langley Research Centre maintains a suite of hypersonic test capabilities spanning Mach~$6$ to Mach~$20$, utilising air, nitrogen, helium, and tetrafluoromethane as test gases~\citep{BergerEtAl2015}. Among these facilities, the Hypersonic Tunnel Facility (HTF)~\citep{ThomasEtAl2010} provides flow conditions representative of Mach~$5$--$7$ flight at altitudes between $20$--$30$~km, generating total temperatures from $1{,}200$--$2{,}200$~K. This facility has supported propulsion studies including the Hypersonic Research Engine programme. The California Institute of Technology T5 shock tunnel~\citep{ParzialeEtAl2013} operates with air or nitrogen at reservoir enthalpies up to $25$~MJ/kg, reservoir pressures of $100$~MPa, and reservoir temperatures reaching $10{,}000$~K. The von K\'{a}rm\'{a}n Institute's Longshot facility tests gases at pressures up to $400$~MPa whilst maintaining moderate temperatures around $2{,}500$~K. Studies on $7$-degree half-angle conical geometries revealed the strong stabilising effect of nosetip bluntness, with wavelet analysis confirming Mack's second mode instabilities~\citep{Grossir2015}. ONERA's F4 facility was originally commissioned for ESA Hermes re-entry studies but has recently been extended for Mach~$8$ flight conditions~\citep{ViguierEtAl2012}. Studies examined real gas influence on flare-induced boundary layer separation, addressing how high-temperature thermochemistry affects shock-boundary layer interactions. CIRA operates the SCIROCCO and GHIBLI plasma wind tunnels capable of temperatures up to $10{,}000$~K, Mach numbers up to $12$, and nozzle exit diameters approaching $1{,}950$~mm and $1{,}150$~mm, respectively~\citep{DeFilippoEtAl2002}. SCIROCCO is the world's largest plasma wind tunnel with arc-jet power up to $70$~MW, whilst GHIBLI employs a nozzle designed for smaller test articles, more recently performing studies on high-enthalpy flows ($h_{0}=20$~MJ/kg, Mach~$8$), achieving excellent flow uniformity in thermochemical non-equilibrium conditions where the method of characteristics cannot be applied~\citep{GuidaSmoraldiSchettino2024}. The Japan Aerospace Exploration Agency (JAXA) operates the High Enthalpy Shock Tunnel (HIEST), a free-piston-driven facility designed to fulfil aerothermodynamic test requirements at hypervelocity conditions of $3$--$7$~km/s~\citep{ItohEtAl2002}. The facility achieves maximum stagnation conditions of $2$~MJ/kg and $150$~MPa, with test durations of approximately $2$~ms. HIEST was constructed to support Japan's HOPE re-entry vehicle and scramjet testing programmes. The University of Oxford operates the T6 Stalker Tunnel, a multi-mode, high-enthalpy, transient ground test facility that combines a free-piston driver with modified barrels from the Oxford Gun Tunnel~\citep{CollenEtAl2021}. Depending on test requirements, T6 can operate as a shock tube, reflected shock tunnel, or expansion tube. The facility achieves enthalpy conditions ranging from $2.7$~MJ/kg to $115$~MJ/kg. A comprehensive review of low- and high-enthalpy facilities can be found in~\cite{Grossir2015}.

\section{Recent Advances in Hypersonic Flow Simulation}\label{section:introduction/recent_advances}

Recent progress in understanding modal instability growth and breakdown to turbulence in hypersonic boundary layers has come from experimental measurements and numerical simulations conducted over the past several decades. Facility-based studies have documented second-mode wave development under various flow conditions, providing detailed measurements of frequency content, growth rates, and spatial evolution of instability waves. Computational work has examined instability mechanisms at resolutions beyond experimental capability, revealing flow structures and energy transfer processes that govern transition. Experimental facilities ranging from quiet wind tunnels to high-enthalpy shock tunnels have enabled investigation across different Mach numbers, Reynolds numbers, and enthalpy levels, whilst numerical approaches from linear stability theory to Direct Numerical Simulation have provided complementary insights into the physics of hypersonic transition. Together, these experimental and computational efforts have established a foundation for understanding how disturbances amplify and eventually break down into turbulent flow in hypersonic boundary layers.


Major hypersonic ground test facilities worldwide have contributed to understanding these phenomena through carefully controlled experiments. The \gls{NASA} Langley Research Centre maintains a suite of hypersonic test capabilities spanning Mach~$6$ to Mach~$20$, utilising air, nitrogen, helium, and tetrafluoromethane as test gases~\citep{BergerEtAl2015}. Among these facilities, the Hypersonic Tunnel Facility (HTF)~\citep{ThomasEtAl2010} provides flow conditions representative of Mach~$5$--$7$ flight at altitudes between $20$--$30$~km, generating total temperatures from $1{,}200$--$2{,}200$~K. This facility has supported propulsion studies including the Hypersonic Research Engine programme. The California Institute of Technology T5 shock tunnel~\citep{ParzialeEtAl2013} operates with air or nitrogen at reservoir enthalpies up to $25$~MJ/kg, reservoir pressures of $100$~MPa, and reservoir temperatures reaching $10{,}000$~K. The von K\'{a}rm\'{a}n Institute's Longshot facility tests gases at pressures up to $400$~MPa whilst maintaining moderate temperatures around $2{,}500$~K. Studies on $7$-degree half-angle conical geometries revealed the strong stabilising effect of nosetip bluntness, with wavelet analysis confirming Mack's second mode instabilities~\citep{Grossir2015}. ONERA's F4 facility was originally commissioned for ESA Hermes re-entry studies but has recently been extended for Mach~$8$ flight conditions~\citep{ViguierEtAl2012}. Studies examined real gas influence on flare-induced boundary layer separation, addressing how high-temperature thermochemistry affects shock-boundary layer interactions. CIRA operates the SCIROCCO and GHIBLI plasma wind tunnels capable of temperatures up to $10{,}000$~K, Mach numbers up to $12$, and nozzle exit diameters approaching $1{,}950$~mm and $1{,}150$~mm, respectively~\citep{DeFilippoEtAl2002}. SCIROCCO is the world's largest plasma wind tunnel with arc-jet power up to $70$~MW, whilst GHIBLI employs a nozzle designed for smaller test articles, more recently performing studies on high-enthalpy flows ($h_{0}=20$~MJ/kg, Mach~$8$), achieving excellent flow uniformity in thermochemical non-equilibrium conditions where the method of characteristics cannot be applied~\citep{GuidaSmoraldiSchettino2024}. The Japan Aerospace Exploration Agency (JAXA) operates the High Enthalpy Shock Tunnel (HIEST), a free-piston-driven facility designed to fulfil aerothermodynamic test requirements at hypervelocity conditions of $3$--$7$~km/s~\citep{ItohEtAl2002}. The facility achieves maximum stagnation conditions of $2$~MJ/kg and $150$~MPa, with test durations of approximately $2$~ms. HIEST was constructed to support Japan's HOPE re-entry vehicle and scramjet testing programmes. The University of Oxford operates the T6 Stalker Tunnel, a multi-mode, high-enthalpy, transient ground test facility that combines a free-piston driver with modified barrels from the Oxford Gun Tunnel~\citep{CollenEtAl2021}. Depending on test requirements, T6 can operate as a shock tube, reflected shock tunnel, or expansion tube. The facility achieves enthalpy conditions ranging from $2.7$~MJ/kg to $115$~MJ/kg. A comprehensive review of low- and high-enthalpy facilities can be found in~\cite{Grossir2015}.


Standard experimental ground-based hypersonic facilities rarely operate at conditions similar to those of a low-altitude high-speed vehicle, as the facility power requirements to replicate the exact combination of velocity, flow density, and therefore stagnation enthalpy can be overwhelming to current electrical grids~\citep{Passiatoreetal2021}. Continuous running and blow down facilities allow the longest running times but are limited to low stagnation enthalpies, whilst reflected shock tunnels and expansion tunnels can achieve high-enthalpy conditions but suffer from short test times~\citep{GuOlivier2020}. No single facility can simultaneously match all flight parameters, as each facility type is capable of matching at least one flow property well whilst sacrificing others. For full-scale propulsion tests at Mach~9, the flow power required can reach 600~MW, demanding driver powers up to 3000~MW, a gap far too large for current mechanical-energy-supported drivers to overcome~\citep{JiangYu2017}.

Given these fundamental limitations in experimental capabilities, high-fidelity computational simulation has become essential for understanding hypersonic flows under conditions that cannot be fully replicated in ground facilities. Proprietary solvers such as \texttt{LAURA} (Langley Aerothermodynamic Upwind Relaxation Algorithm)~\citep{Cheatwood1996} and \texttt{DPLR} (Data-Parallel Line Relaxation)~\citep{Candler1994} have been extensively validated against flight data and ground test measurements for atmospheric entry vehicles, incorporating thermochemical nonequilibrium models, finite-rate chemistry, and surface ablation capabilities. These codes have provided critical aerothermal predictions for missions including Space Shuttle, Mars Science Laboratory, and Orion. \texttt{US3D}~\citep{Nompelis2004} and \texttt{LeMANS}~\citep{Scalabrin2005} offer comprehensive frameworks for hypersonic simulation on unstructured grids with capabilities ranging from Reynolds-Averaged Navier-Stokes to Direct Numerical Simulation of transition phenomena. \texttt{DLR-TAU}~\citep{Mack2002} extends aircraft CFD capabilities to hypersonic regimes with high-temperature gas models and has been applied to European space programmes.


Open-source solvers have gained prominence through community-driven development and accessible implementations. \texttt{Eilmer}~\citep{Gollan2013} is a compressible flow simulation code with particular emphasis on high-temperature gas dynamics, featuring flexible gas models ranging from calorically perfect to finite-rate chemistry and shock-fitting capabilities with extensive validation against shock tunnel experiments. The solver's Python-based interface enables scriptable workflow automation and integration with optimisation frameworks. \texttt{SU2-NEMO}~\citep{MaierEtAl2021} extends the widely-used SU2 framework to hypersonic nonequilibrium flows, uniquely integrating adjoint capabilities for sensitivity analysis and design optimisation of hypersonic configurations. The code implements comprehensive thermochemical modelling including multiple temperature models, detailed multicomponent transport properties, and flexible chemistry mechanisms. \texttt{hy2Foam}~\citep{Casseau2016a,Casseau2016b} leverages the OpenFOAM framework with hybrid DSMC-CFD capabilities for cross-regime simulations from continuum to rarefied flows, whilst its integration with OpenFOAM's conjugate heat transfer framework enables coupled fluid-structure-thermal simulations relevant to thermal protection system design. The \texttt{HTR solver}~\citep{DiRenzoFuUrzay2020} represents a modern approach employing task-based parallelism through the Legion programming model and high-order entropy-stable finite-difference methods optimised for modern GPU architectures, enabling Direct Numerical Simulation and Large Eddy Simulation at unprecedented Reynolds numbers that bridge the gap between laboratory and flight conditions.

Accurate thermochemical modelling requires specialised libraries for gas properties under extreme conditions. \texttt{CEA} (Chemical Equilibrium with Applications)~\citep{McBrideGordon1996} computes chemical equilibrium compositions and thermodynamic properties for complex mixtures through Gibbs free energy minimisation, including an extensive thermodynamic database covering hundreds of species relevant to combustion and atmospheric entry. Whilst limited to equilibrium assumptions, CEA serves as a fundamental tool for initial condition specification and validating more complex kinetic models. \texttt{Cantera}~\citep{GoodwinMoffatSpeth2015} provides an open-source framework for chemical kinetics, thermodynamics, and transport processes with finite-rate chemistry using arbitrary reaction mechanisms. The library computes thermodynamic properties using NASA polynomial fits and transport properties using kinetic theory with collision integrals, offering Python, MATLAB, and C++ interfaces that facilitate integration with various flow solvers. \texttt{KAPPA} (Kinetics And Plasma Properties Analyser)~\citep{CampoliEtAl2019} specialises in high-temperature atmospheric entry with plasma modelling capabilities including ionisation kinetics and electron-impact processes, essential for high-speed Earth entries where ionisation significantly affects flowfield structure and radiative heating. The library implements state-of-the-art transport property models with collision integrals from ab initio calculations. \texttt{Mutation++} (MUlticomponent Thermodynamic And Transport properties for IONised gases in C++)~\citep{Scoggins2020} offers comprehensive transport properties and chemical production rates for multicomponent mixtures, implementing multiple mixture models including equilibrium, thermal nonequilibrium, and state-specific kinetics approaches. The library features advanced Chapman-Enskog transport calculations using various approximations suitable for different flow regimes and computational cost constraints, with extensive validation against experimental data from plasma wind tunnels and shock tubes ensuring accuracy under extreme conditions.



The development of sophisticated computational tools and specialised thermochemical libraries over the past two decades has significantly expanded hypersonic research capabilities. High-fidelity simulations now access parameter spaces beyond ground facility limits, resolve instability mechanisms at scales too fine for experimental measurements, and enable systematic studies of thermochemical effects on transition and turbulence. Experiments remain essential for validation and investigating phenomena driven by facility-specific disturbance environments, but the combination of advanced computational frameworks with targeted experimental campaigns has produced major advances in understanding hypersonic flow physics. With continued improvements in high-performance computing, numerical simulation has become an essential tool for hypersonic vehicle development and fundamental research into extreme flow conditions.

% Open-source solvers have gained prominence through community-driven development and accessible implementations. \texttt{Eilmer}~\citep{Gollan2013} features flexible gas models and shock-fitting capabilities with extensive validation against shock tunnel experiments. \texttt{SU2-NEMO}~\citep{MaierEtAl2021} uniquely integrates adjoint capabilities for sensitivity analysis and design optimisation of hypersonic configurations. \texttt{hy2Foam}~\citep{Casseau2016a,Casseau2016b} leverages the OpenFOAM framework with hybrid DSMC-CFD capabilities for cross-regime simulations. The \texttt{HTR solver}~\citep{DiRenzoFuUrzay2020} employs task-based parallelism and high-order methods optimised for modern GPU architectures, enabling simulations at unprecedented Reynolds numbers.

% Accurate thermochemical modelling requires specialised libraries for gas properties under extreme conditions. \texttt{CEA}~\citep{McBrideGordon1996} computes chemical equilibrium compositions through Gibbs free energy minimisation, whilst \texttt{Cantera}~\citep{GoodwinMoffatSpeth2015} provides open-source finite-rate chemistry with arbitrary reaction mechanisms. \texttt{KAPPA}~\citep{CampoliEtAl2019} specialises in high-temperature atmospheric entry with plasma modelling and ionisation kinetics. \texttt{Mutation++}~\citep{Scoggins2020} offers comprehensive transport properties and chemical production rates for multicomponent mixtures, implementing state-specific kinetics and advanced Chapman-Enskog transport calculations validated against plasma wind tunnel data.

% Given these fundamental limitations in experimental capabilities, high-fidelity computational simulation has become essential for understanding hypersonic flows under conditions that cannot be fully replicated in ground facilities. The computational simulation of hypersonic flows requires sophisticated numerical solvers capable of handling the complex physical phenomena that arise at extreme velocities and temperatures. Various CFD solvers have been developed specifically for nonequilibrium flows in the near-continuum regime, each with distinct capabilities and modelling approaches. \texttt{LAURA} (Langley Aerothermodynamic Upwind Relaxation Algorithm)~\citep{Cheatwood1996} represents one of the pioneering efforts in hypersonic CFD development. This three-dimensional, finite-volume solver is designed for computing hypersonic flowfields in thermochemical nonequilibrium. The code employs an upwind point-implicit relaxation scheme for solving the Navier-Stokes equations on structured grids. LAURA has been extensively validated against flight data and ground test measurements, particularly for atmospheric entry vehicles. Its capabilities include finite-rate chemistry modelling with multiple species, thermal nonequilibrium with separate translational and vibrational temperatures, and surface catalysis effects. The solver has been applied to numerous missions including Space Shuttle, Mars Science Laboratory, and Orion crew vehicle, providing critical aerothermal predictions for heat shield design. LAURA's robust performance on entry configurations with strong shock interactions and separated flow regions has made it a workhorse code for entry vehicle programmes.



% The computational simulation of hypersonic flows requires sophisticated numerical solvers capable of handling the complex physical phenomena that arise at extreme velocities and temperatures. Various CFD solvers have been developed specifically for nonequilibrium flows in the near-continuum regime, each with distinct capabilities and modeling approaches. \texttt{LAURA} (Langley Aerothermodynamic Upwind Relaxation Algorithm)~\citep{Cheatwood1996} represents one of the pioneering efforts in hypersonic CFD development. This three-dimensional, finite-volume solver is designed for computing hypersonic flowfields in thermochemical nonequilibrium. The code employs an upwind point-implicit relaxation scheme for solving the Navier-Stokes equations on structured grids. LAURA has been extensively validated against flight data and ground test measurements, particularly for atmospheric entry vehicles. Its capabilities include finite-rate chemistry modeling with multiple species, thermal nonequilibrium with separate translational and vibrational temperatures, and surface catalysis effects. The solver has been applied to numerous missions including Space Shuttle, Mars Science Laboratory, and Orion crew vehicle, providing critical aerothermal predictions for heat shield design. LAURA's robust performance on entry configurations with strong shock interactions and separated flow regions has made it a workhorse code for entry vehicle programs.

% \texttt{DPLR} (Data-Parallel Line Relaxation)~\citep{Candler1994} is another prominent hypersonic flow solver optimized for massively parallel computing architectures. DPLR solves the reacting Navier-Stokes equations using finite-volume discretization on structured grids with implicit time integration. The code implements detailed thermochemical models including Park's two-temperature model for thermal nonequilibrium and multiple chemical kinetics mechanisms. DPLR's parallel scalability enables high-fidelity simulations of complete entry vehicle configurations at flight conditions. The solver incorporates advanced physical models for surface ablation, allowing coupled simulations where the heat shield material response feeds back into the flowfield solution. DPLR has been extensively used for Mars entry vehicle design, including Mars Pathfinder, Mars Exploration Rovers, and Mars 2020 Perseverance, as well as for lunar and Earth return missions. Its radiation modeling capabilities, essential for high-speed Earth entries, include tangent-slab and ray-tracing methods for computing radiative heating from shock-layer emission.

% \texttt{LeMANS} (Michigan Aerothermodynamic Navier-Stokes)~\citep{Scalabrin2005} is a parallel three-dimensional Navier-Stokes solver for hypersonic reacting flows. LeMANS employs finite-volume discretization on unstructured grids, providing geometric flexibility for complex configurations. The code implements modified Steger-Warming flux vector splitting for inviscid fluxes and central differencing for viscous terms. LeMANS features comprehensive thermochemical modeling including Park's two-temperature model, multiple chemical kinetics mechanisms, and various transport property formulations. The solver's modular architecture facilitates implementation of new physical models, making it popular in the academic community for hypersonic research. LeMANS has been applied to fundamental transition studies, scramjet combustor simulations, and entry vehicle aerothermodynamics. Its unstructured grid capability enables automated mesh adaptation in regions of strong gradients, improving solution accuracy while managing computational cost.

% \texttt{US3D} (Unstructured 3D)~\citep{Nompelis2004} represents a comprehensive framework for hypersonic flow simulation on unstructured grids. US3D solves the compressible Navier-Stokes equations with finite-rate chemistry and thermal nonequilibrium using cell-centered finite-volume methods. The code implements multiple flux reconstruction schemes including AUSM+-up and Roe with various limiters for shock capturing. US3D's thermochemical modeling includes multiple temperature models (one, two, or multiple temperature) and flexible chemical kinetics with user-specified reaction mechanisms. The solver features advanced turbulence modeling capabilities including Reynolds-Averaged Navier-Stokes and Large Eddy Simulation approaches adapted for hypersonic flows. US3D has been extensively validated for hypersonic boundary layer transition, demonstrating excellent agreement with quiet wind tunnel experiments. The code's high-order spatial accuracy options enable Direct Numerical Simulation of transition phenomena, providing detailed insights into instability growth and breakdown mechanisms not accessible through experiments.

% \texttt{DLR-TAU}~\citep{Mack2002} is a versatile CFD package originally designed for aircraft applications but extended to hypersonic flow regimes. TAU employs unstructured hybrid grids (tetrahedra, prisms, pyramids, hexahedra) with cell-centered and vertex-centered finite-volume formulations. For hypersonic applications, TAU incorporates high-temperature gas models with chemical reactions, thermal nonequilibrium, and radiative heat transfer. The code implements multiple turbulence models including Spalart-Allmaras, $k-\omega$, and Reynolds Stress Models, with adaptations for compressibility effects. TAU's comprehensive pre- and post-processing tools facilitate rapid configuration setup and analysis. The solver has been applied to European space programs including the Intermediate eXperimental Vehicle (IXV) and Space Rider, providing aerothermodynamic databases for design. TAU's adaptive mesh refinement capabilities enable efficient resolution of localized flow features such as shock-boundary layer interactions.

% \texttt{Eilmer}~\citep{Gollan2013} is a compressible flow simulation code with particular emphasis on high-temperature gas dynamics. Eilmer solves the Euler and Navier-Stokes equations on structured and unstructured grids using finite-volume methods. The code features a flexible gas model library supporting calorically perfect gas, thermally perfect gas, chemical equilibrium, and finite-rate chemistry models. Eilmer implements various flux calculators including Riemann solvers (Roe, HLLE, AUSMDV) suitable for hypersonic flows with strong shocks. The solver's Python-based interface enables scriptable workflow automation and integration with optimization frameworks. Eilmer has been extensively validated against shock tunnel experiments, particularly for several high-enthalpy facilities. The code's shock-fitting capabilities, in addition to shock-capturing, provide highly accurate shock resolution for fundamental studies. Eilmer's open-source availability and active development community have made it popular for academic hypersonic research.

% \texttt{SU2-NEMO} (Nonequilibrium flow with MUlticomponent transport)~\citep{MaierEtAl2021} extends the widely-used SU2 framework to hypersonic nonequilibrium flows. SU2-NEMO implements finite-volume discretization on unstructured grids with edge-based data structures for computational efficiency. The solver features comprehensive thermochemical modeling including multiple temperature models, detailed transport properties with multicomponent diffusion, and flexible chemistry mechanisms. SU2-NEMO's integration with SU2's adjoint capabilities enables sensitivity analysis and design optimization for hypersonic configurations, a unique capability among hypersonic solvers. The code implements multiple flux schemes including AUSM-based and Roe-type methods with entropy fixes. SU2-NEMO's open-source nature and active development community have driven rapid feature expansion, including recent additions for turbulence modeling and conjugate heat transfer. The solver has been validated against fundamental shock-tube experiments and applied to entry vehicle configurations including Stardust and Fire II.

% The \texttt{HTR solver} (Hypersonic Task-based Research)~\citep{DiRenzoFuUrzay2020} represents a modern approach leveraging task-based parallelism for extreme-scale computing. HTR implements high-order finite-difference methods on structured grids for direct numerical simulation and large eddy simulation of hypersonic turbulent flows. The solver employs the Legion programming model, enabling efficient execution on heterogeneous architectures including GPUs. HTR's high-order spatial discretization with entropy-stable schemes provides accurate shock capturing without excessive dissipation in smooth regions. The code features detailed thermochemical modeling including mixture-averaged and multicomponent transport, Park's two-temperature model, and finite-rate chemistry. HTR has been applied to fundamental studies of shock-turbulence interaction, hypersonic boundary layer transition, and scramjet combustion. The solver's ability to exploit modern supercomputing architectures enables simulations at unprecedented Reynolds numbers, bridging the gap between laboratory and flight conditions.

% \texttt{hy2Foam}~\citep{Casseau2016a,Casseau2016b} is an open-source hypersonic flow solver built on the OpenFOAM framework. hy2Foam solves the compressible Navier-Stokes equations with chemical reactions and thermal nonequilibrium using finite-volume discretization on arbitrary polyhedral grids. The solver implements multiple thermochemical models including Park's two-temperature model, vibration-dissociation coupling, and various chemical kinetics mechanisms. hy2Foam leverages OpenFOAM's extensive utility libraries for mesh generation, manipulation, and post-processing. The code features hybrid DSMC-CFD capabilities, enabling simulation across flow regimes from continuum to rarefied. hy2Foam has been validated against experimental data from hypersonic wind tunnels and applied to reentry vehicle configurations and high-altitude propulsion systems. Its integration with OpenFOAM's conjugate heat transfer framework enables coupled fluid-structure-thermal simulations relevant to thermal protection system design.

% Complementing these flow solvers are specialized thermochemical libraries that provide accurate gas properties essential for hypersonic simulations. \texttt{CEA} (Chemical Equilibrium with Applications)~\citep{McBrideGordon1996} computes chemical equilibrium compositions and thermodynamic properties for complex mixtures. CEA implements Gibbs free energy minimization to determine equilibrium compositions for specified temperature, pressure, and elemental constraints. The program includes an extensive thermodynamic database covering hundreds of species relevant to combustion and atmospheric entry. CEA provides properties including enthalpy, entropy, specific heat, density, and speed of sound for equilibrium mixtures. While limited to equilibrium assumptions, CEA serves as a fundamental tool for initial condition specification, understanding limiting behavior, and validating more complex kinetic models.

% \texttt{Cantera}~\citep{GoodwinMoffatSpeth2015} provides an open-source framework for chemical kinetics, thermodynamics, and transport processes. Cantera implements detailed chemical kinetics with arbitrary reaction mechanisms, enabling finite-rate chemistry simulations. The library computes thermodynamic properties using NASA polynomial fits and transport properties using kinetic theory with collision integrals. Cantera's Python, MATLAB, and C++ interfaces facilitate integration with various flow solvers and enable stand-alone chemical kinetics analysis. The framework supports multiple reactor models (zero-dimensional, one-dimensional flames, plug flow) useful for detailed chemistry studies. Cantera has been widely adopted in combustion research and increasingly in hypersonic applications requiring detailed chemistry modeling.

% \texttt{KAPPA} (Kinetics And Plasma Properties Analyzer)~\citep{CampoliEtAl2019} is a specialized library for high-temperature gas properties relevant to atmospheric entry. KAPPA computes thermodynamic properties, transport coefficients, and chemical kinetics for mixtures in thermal and chemical nonequilibrium. The library implements state-of-the-art models for transport properties including collision integrals from ab initio calculations and semi-empirical correlations. KAPPA features plasma modeling capabilities including ionization kinetics and electron-impact processes, essential for high-speed Earth entries where ionization significantly affects flowfield structure and radiative heating. The library's modular architecture enables coupling with various CFD solvers.

% \texttt{Mutation++} (MUlticomponent Thermodynamic And Transport properties for IONized gases in C++)~\citep{Scoggins2020} is a comprehensive library for high-temperature gas mixture properties. Mutation++ provides thermodynamic properties, transport coefficients, chemical production rates, and energy exchange terms for multicomponent mixtures. The library implements multiple mixture models including equilibrium, thermal nonequilibrium (two-temperature), and state-specific kinetics approaches. Mutation++ features advanced transport property calculations using various approximations (Chapman-Enskog, effective binary diffusion, Ramshaw method) suitable for different flow regimes and computational cost constraints. The library includes an extensive database of species thermodynamic data and collision parameters. Mutation++'s C++ implementation with Python bindings facilitates integration with modern flow solvers. The library has been extensively validated against experimental data from various plasma wind tunnels and shock tubes, ensuring accuracy under extreme conditions. Mutation++'s active development continues to incorporate latest advances in high-temperature gas physics, including improved electronic excitation models and radiation property calculations.


% Standard experimental ground-based hypersonic facilities rarely operate at conditions similar to those of a low-altitude high-speed vehicle, as the facility power requirements to replicate the exact combination of velocity, flow density, and therefore stagnation enthalpy can be overwhelming to current electrical grids~\citep{Passiatoreetal2021}. Continuous running and blow down facilities allow the longest running times but are limited to low stagnation enthalpies, whilst reflected shock tunnels and expansion tunnels can achieve high-enthalpy conditions but suffer from short test times~\citep{GuOlivier2020}. No single facility can simultaneously match all flight parameters, as each facility type is capable of matching at least one flow property well whilst sacrificing others. For full-scale propulsion tests at Mach 9, the flow power required can reach 600 MW, demanding driver powers up to 3000 MW, a gap far too large for current mechanical-energy-supported drivers to overcome  ~\citep{JiangYu2017}.


% One of the early studies to understand the transition mechanism in hypersonic boundary layers was conducted by Parziale et al.~\citep{Parziale2015} in the Caltech T5 hypervelocity shock tunnel, where differential interferometry was used to measure instability waves under high-enthalpy conditions. This work, along with subsequent investigations in facilities such as the Boeing/AFOSR Mach-6 Quiet Tunnel~\citep{Hofferth2013, Tang2015} and the DLR HEG shock tunnel~\citep{Laurence2016}, established that second-mode instability dominates the transition process at hypersonic Mach numbers. These experimental efforts revealed that high-enthalpy effects, including thermochemical nonequilibrium and vibrational excitation, significantly influence the characteristics and growth rates of second-mode waves~\citep{Salemi2018, Hameed2024}.


% BergerEtAl2015

% BergerEtAl2015

% The \gls{NASA} Langley Research Centre maintains a suite of hypersonic test capabilities spanning Mach~$6$ to Mach~$20$, utilising air, nitrogen, helium, and tetrafluoromethane as test gases~\citep{BergerEtAl2015}. Among these facilities, the Hypersonic Tunnel Facility (HTF)~\citep{ThomasEtAl2010} provides flow conditions representative of Mach~$5$--$7$ flight at altitudes between $20$--$30$~km, generating total temperatures from $1200$--$2200$~K. This facility has supported propulsion studies including the Hypersonic Research Engine programme. The California Institute of Technology T5 shock tunnel~\citep{ParzialeEtAl2013} operates with air or nitrogen at reservoir enthalpies up to $25$~MJ/kg, reservoir pressures of $100$~MPa, and reservoir temperatures reaching $10{,}000$~K. The von Karman Institute's \texttt{Longshot facility} tests gases at pressures up to 400 MPa while maintaining moderate temperatures around 2500 K. Studies on 7-degree half-angle conical geometries revealed the strong stabilizing effect of nosetip bluntness, with wavelet analysis confirming Mack's second mode instabilities~\citep{Grossir2015}.ONERA \texttt{F4 facility} was originally commissioned for ESA Hermes re-entry studies, but has recently been extended for Mach~8 flight conditions~\citep{ViguierEtAl2012}. Studies examined real gas influence on flare-induced boundary layer separation, addressing how high-temperature thermochemistry affects shock-boundary layer interactions. The \texttt{CIRA} operated  \texttt{SCIROCCO} and \texttt{GHIBLI} plasma tunnels capable of temperatures up to $10,000$ K, Mach numbers up to $12$, and nozzle exit diameters approaching $1950$~mm and $1150$~mm each~\citep{DeFilippoEtAl2002}. \texttt{SCIROCCO} is the world's largest plasma wind tunnel with arc-jet power up to 70 MW, while \texttt{GHIBLI} employs a nozzle designed for smaller cases, more recently performing studies on high-enthalpy flows ($h_{0}=20$ MJ/kg, Mach 8), achieving excellent flow uniformity in thermochemical non-equilibrium conditions where the method of characteristics cannot be applied~\cite{GuidaSmoraldiSchettino2024}.

% The \gls{NASA} Langley Research Centre maintains a suite of hypersonic test capabilities spanning Mach~$6$ to Mach~$20$, utilising air, nitrogen, helium, and tetrafluoromethane as test gases~\citep{BergerEtAl2015}. Among these facilities, the Hypersonic Tunnel Facility (HTF)~\citep{ThomasEtAl2010} provides flow conditions representative of Mach~$5$--$7$ flight at altitudes between $20$--$30$~km, generating total temperatures from $1{,}200$--$2{,}200$~K. This facility has supported propulsion studies including the Hypersonic Research Engine programme. The California Institute of Technology T5 shock tunnel~\citep{ParzialeEtAl2013} operates with air or nitrogen at reservoir enthalpies up to $25$~MJ/kg, reservoir pressures of $100$~MPa, and reservoir temperatures reaching $10{,}000$~K. The von K\'{a}rm\'{a}n Institute's Longshot facility tests gases at pressures up to $400$~MPa while maintaining moderate temperatures around $2{,}500$~K. Studies on $7$-degree half-angle conical geometries revealed the strong stabilising effect of nosetip bluntness, with wavelet analysis confirming Mack's second mode instabilities~\citep{Grossir2015}. ONERA's F4 facility was originally commissioned for ESA Hermes re-entry studies, but has recently been extended for Mach~$8$ flight conditions~\citep{ViguierEtAl2012}. Studies examined real gas influence on flare-induced boundary layer separation, addressing how high-temperature thermochemistry affects shock-boundary layer interactions. CIRA operates the SCIROCCO and GHIBLI plasma wind tunnels capable of temperatures up to $10{,}000$~K, Mach numbers up to $12$, and nozzle exit diameters approaching $1{,}950$~mm and $1{,}150$~mm, respectively~\citep{DeFilippoEtAl2002}. SCIROCCO is the world's largest plasma wind tunnel with arc-jet power up to $70$~MW, while GHIBLI employs a nozzle designed for smaller test articles, more recently performing studies on high-enthalpy flows ($h_{0}=20$~MJ/kg, Mach~$8$), achieving excellent flow uniformity in thermochemical non-equilibrium conditions where the method of characteristics cannot be applied~\citep{GuidaSmoraldiSchettino2024}. The Japan Aerospace Exploration Agency (JAXA) operates the High Enthalpy Shock Tunnel (HIEST), a free-piston-driven facility designed to fulfil aerothermodynamic test requirements at hypervelocity conditions of $3$--$7$~km/s~\citep{ItohEtAl2002}. The facility achieves maximum stagnation conditions of $2$~MJ/kg and $150$~MPa, with test durations of approximately $2$~ms. HIEST was constructed to support Japan's HOPE re-entry vehicle and scramjet testing programmes. The University of Oxford operates the T6 Stalker Tunnel, a multi-mode, high-enthalpy, transient ground test facility that combines a free-piston driver with modified barrels from the Oxford Gun Tunnel~\citep{CollenEtAl2021}. Depending on test requirements, T6 can operate as a shock tube, reflected shock tunnel, or expansion tube. The facility achieves enthalpy conditions ranging from $2.7$~MJ/kg to $115$~MJ/kg. A comprehensive literature about low and high enthalpy facilities can be found in \cite{Grossir2015}

% Oxford~\citep{CollenEtAl2021}


% The Japan Aerospace Exploration Agency (JAXA) operates the High Enthalpy Shock Tunnel (HIEST), a free-piston-driven facility designed to fulfil aerothermodynamic test requirements at hypervelocity conditions of $3$--$7$~km/s~\citep{ItohEtAl1999}. The facility was constructed at the National Aerospace Laboratory (NAL) Kakuda facility to support Japan's HOPE re-entry vehicle and scramjet testing programmes. HIEST consists of a compression tube ($16$~m long, $600$~mm diameter), a shock tube ($17$~m long, $180$~mm diameter), and employs five interchangeable pistons ranging from $220$--$780$~kg to achieve the desired stagnation enthalpy and pressure conditions. For aerodynamic tests, a conical nozzle with throat diameter $24$--$50$~mm and exit diameter $1{,}200$~mm enables testing of large-span models, whilst scramjet tests utilise a contoured nozzle with throat and exit diameters of $50$~mm and $800$~mm, respectively. The facility achieves maximum stagnation conditions of $2$~MJ/kg and $150$~MPa, with test durations of approximately $2$~ms at the highest enthalpy conditions.



% The \gls{NASA} Langley Research Center operates tunnels that can achieve Mach~6-Mach~20 flight with various test gases, air, pure nitrogen, helium, and tetrafluoromethane (CF4). Another high-enthalpy facility at \gls{NASA} achieves conditions up to Mach 7, which is predominantly used for rockets. The von Karman Institute's \texttt{Longshot facility} tests gases at pressures up to 400 MPa while maintaining moderate temperatures around 2500 K. Studies on 7-degree half-angle conical geometries revealed the strong stabilizing effect of nosetip bluntness, with wavelet analysis confirming Mack's second mode instabilities~\citep{Grossir2015}. ONERA \texttt{F4 facility} was originally commissioned for ESA Hermes re-entry studies, but has recently been extended for Mach~8 flight conditions~\citep{ViguierEtAl2012}. Studies examined real gas influence on flare-induced boundary layer separation, addressing how high-temperature thermochemistry affects shock-boundary layer interactions. The \texttt{CIRA} operated  \texttt{SCIROCCO} and \texttt{GHIBLI} plasma tunnels capable of temperatures up to $10,000$ K, Mach numbers up to $12$, and nozzle exit diameters approaching $1950$~mm and $1150$~mm each~\citep{DeFilippoEtAl2002}. \texttt{SCIROCCO} is the world's largest plasma wind tunnel with arc-jet power up to 70 MW, while \texttt{GHIBLI} employs a nozzle designed for smaller cases, more recently performing studies on high-enthalpy flows ($h_{0}=20$ MJ/kg, Mach 8), achieving excellent flow uniformity in thermochemical non-equilibrium conditions where the method of characteristics cannot be applied~\cite{GuidaSmoraldiSchettino2024}. \texttt{HIEST} operates across enthalpies from 3 to 20 MJ/kg with stagnation pressures between 10 and 60 MPa~\cite{hiest_wind_tunnel_test, hiest_apollo}. Apollo capsule testing addressed boundary layer trip effectiveness with trip heights between 0.2 and 1.3 mm~\cite{hiest_apollo}, while HTV-R capsule aerodynamics testing examined trim angle characteristics under high-temperature real-gas effects~\cite{hiest_free_flight}. Fundamental boundary layer transition research employs 1100 mm-length 7-degree half-angle cones examining surface pressure fluctuation growth in chemically reacting flows where oxygen dissociation reaches approximately 40\%~\cite{hiest_research_works}. Transpiration cooling investigations on flat plate geometries examined coolant injection effects for active thermal management~\cite{hiest_research_works}.  The Oxford \texttt{T6 Stalker Tunnel}, Europe's fastest wind tunnel, produces flows exceeding 20 km/s with maximum total pressures of 75 MPa and temperatures up to 5,000 K~\cite{t6_nwtf}. Multi-mode operation includes reflected shock tunnel, expansion tube, and shock tube configurations~\cite{t6_development}. HIFiRE-5 model testing examined crossflow instability and stationary wave amplification on elliptic cone geometries~\cite{t6_afrl}, while hypersonic intake studies including CREST intake configurations at Mach 5~\cite{t6_afrl} addressed inlet-isolator unstart phenomena. Purdue University's \texttt{quiet tunnel} enables low-disturbance boundary layer transition studies. China's \texttt{JF-22}, spanning 167 meters with a 4-meter diameter, simulates altitudes from 40 to 90 kilometers at speeds up to 10 km/s (approximately Mach 30)~\cite{jf22_unveiling, jf22_chinadaily}. The facility employs precisely timed explosions generating converging shock waves rather than conventional compression~\cite{jf22_newatlas} to study gliding aircraft, space-to-Earth flight vehicles, and multistage launch vehicles.

% Despite these extensive experimental capabilities, the computational cost to perform full-scale, high-fidelity simulations of hypersonic glide vehicles remains prohibitive with current technology. Direct numerical simulation of realistic flight conditions would require resolving turbulent boundary layers, thermochemical nonequilibrium, and multiscale physics across entire vehicle geometries—computational demands that exceed present supercomputing resources by orders of magnitude.



% Despite these extensive experimental capabilities, fundamental challenges persist in both computational and experimental hypersonic research. Standard ground-based hypersonic facilities rarely operate at conditions similar to those of low-altitude high-speed vehicles, as the facility power requirements to replicate the exact combination of velocity, flow density, and stagnation enthalpy can overwhelm current electrical grids~\citep{passiatore2022}. Meanwhile, the computational cost to perform full-scale, high-fidelity simulations of hypersonic glide vehicles remains prohibitive with current technology. Direct numerical simulation of realistic flight conditions would require resolving turbulent boundary layers, thermochemical nonequilibrium, and multiscale physics across entire vehicle geometries—computational demands that exceed present supercomputing resources by orders of magnitude. Laminar and turbulent flows have received significant attention at sub-orbital enthalpies, albeit with approximations about recombination and dissociation reaction mechanisms. Several studies focusing on simplified turbulent hypersonic boundary layers with minimal chemical reactions were developed~\cite{martin1998}, though rigorous differentiation between low-enthalpy and high-enthalpy calculations was not provided until later work~\cite{duanmartin2011} for gases with little chemical activity while neglecting thermal non-equilibrium. These studies have mostly considered temporal evolution of the boundary layer with minimal chemical activity, leaving the nature of flow physics at extreme hypersonic flight conditions relatively unexplored~\citep{Candler2019}.


% The \gls{NASA} \texttt{Langley Research Center} operates numerous wind tunnels ranging from Mach~6 to Mach~20 randing from air, pure Nitrogen, Helium, and Tetrafluoromethane, $\ce{CF4}$. Another high-enthalpy facility at NASA and go up to Mach 7 conditions. The von Karman Institute's \texttt{Longshot facility} test gases to pressures up to 400 MPa while maintaining moderate temperatures around 2500 K~\cite{vki_longshot_capabilities},Studies on 7-degree half-angle conical geometries revealed the strong stabilizing effect of nosetip bluntness, with wavelet analysis confirming Mack's second mode instabilities~\citep{Grossir2015}.  Italy's \texttt{CIRA} operates \texttt{SCIROCCO}, the world's largest plasma wind tunnel for testing very large models in thermal protection system qualification~\cite{scirocco_characterization}. The plasma jet reaches temperatures up to 10,000 K, Mach numbers up to 10, and nozzle exit diameters approaching 2 meters~\cite{scirocco_numerical_assessment}, with interchangeable conical nozzles of 1150 mm and 1950 mm diameter~\cite{scirocco_characterization}. Studies examined real gas influence on flare-induced boundary layer separation, addressing how high-temperature thermochemistry affects shock-boundary layer interactions~\cite{f4_validation}. The facility was extensively developed for scramjet combustion testing, with a 1.5-meter LAPCAT II small-scale flight experiment configuration tested at Mach 7.4 conditions, demonstrating positive aero-propulsive balance~\cite{f4_scramjet_tests}. Japan's HIEST operates across enthalpies from 3 to 20 MJ/kg with stagnation pressures between 10 and 60 MPa~\cite{hiest_wind_tunnel_test, hiest_apollo}. For capsule studies, Apollo capsule testing addressed boundary layer trip effectiveness with trip heights between 0.2 and 1.3 mm~\cite{hiest_apollo}, while HTV-R capsule aerodynamics testing examined trim angle characteristics under high-temperature real-gas effects~\cite{hiest_free_flight}. Separately, fundamental boundary layer transition research employs slender cone geometries, with studies on 1100mm-length 7-degree half-angle cones examining surface pressure fluctuation growth in chemically reacting flows where oxygen dissociation reaches approximately 40\%~\cite{hiest_research_works}. Transpiration cooling investigations on flat plate geometries examined coolant injection effects~\cite{hiest_research_works}, exploring active thermal management for sustained hypersonic cruise vehicles. The Oxford T6 Stalker Tunnel, Europe's fastest wind tunnel, produces flows exceeding 20 km/s with maximum total pressures of 75 MPa and temperatures up to 5,000 K~\cite{t6_nwtf}. Multi-mode operation includes Reflected Shock Tunnel, Expansion Tube, and shock tube configurations~\cite{t6_development}. Research encompasses shock layer thermochemistry and radiation, convective heating, and boundary layer transition~\cite{t6_development}. HIFiRE-5 model testing examined crossflow instability and stationary wave amplification~\cite{t6_afrl} on elliptic cone geometries, representing fundamental transition research on three-dimensional boundary layers, while hypersonic intake studies including CREST intake configurations at Mach 5~\cite{t6_afrl} addressed inlet-isolator unstart phenomena under propulsion-relevant conditions. China's JF-22, spanning 167 meters in length with a 4-meter diameter, simulates altitudes from 40 to 90 kilometers at speeds up to 10 km/s---approximately Mach 30~\cite{jf22_unveiling, jf22_chinadaily}. JF-22 employs precisely timed explosions generating converging shock waves rather than conventional compression~\cite{jf22_newatlas}. Research includes gliding aircraft, space-to-Earth flight vehicles, and multistage launch vehicle




% Hypersonic wind tunnels represent critical infrastructure for advancing aerospace technology, enabling ground-based testing of vehicles designed to fly at speeds exceeding Mach 5. These facilities vary significantly in their operational principles, test conditions, and research focus, collectively spanning the entire hypersonic flight regime from atmospheric entry to scramjet propulsion testing.



% The von Karman Institute's Longshot facility in Belgium compresses test gases to pressures up to 400 MPa while maintaining moderate temperatures around 2500 K~\cite{vki_longshot_capabilities}, with free-stream Mach numbers ranging between 9.5 and 12~\cite{grossir_space_debris}. Research emphasizes low-enthalpy, high-Reynolds-number conditions that approach perfect gas behavior. Studies on 7-degree half-angle conical geometries revealed the strong stabilizing effect of nosetip bluntness, with wavelet analysis confirming Mack's second mode instabilities~\cite{grossir_transition_studies}. 




% The von Karman Institute's Longshot facility in Belgium stands as one of the pioneering gun tunnels for hypersonic research. This facility compresses test gases to pressures up to 400 MPa while maintaining moderate temperatures around 2500 K~\cite{vki_longshot_capabilities}, with free-stream Mach numbers ranging between 9.5 and 12 and unit Reynolds numbers within 3--10 $\times$ 10$^{6}$/m~\cite{grossir_space_debris}. The Longshot's research emphasizes fundamental boundary layer transition phenomena on canonical geometries. Studies on 7-degree half-angle conical geometries with varying nosetip radii revealed the strong stabilizing effect of nosetip bluntness, with wavelet analysis confirming Mack's second mode instabilities~\cite{grossir_transition_studies}. These cone experiments at high Reynolds numbers specifically examined disturbance evolution from linear growth through nonlinear breakdown to fully turbulent flow, with LIF-based Schlieren visualization revealing rope-shaped structures. Separately, the facility has conducted free-flight measurement techniques for space debris aerodynamic characterization on hemispherical and annular geometries~\cite{grossir_space_debris}, representing blunt-body capsule configurations, while approximately 140 tests with carbon dioxide enabled Martian entry simulations~\cite{ilich_martian_entries}.

% Italy's CIRA operates SCIROCCO, the world's largest plasma wind tunnel for testing very large models in thermal protection system qualification~\cite{scirocco_characterization}. The plasma jet reaches temperatures up to 10,000 K, Mach numbers up to 10, and nozzle exit diameters approaching 2 meters~\cite{scirocco_numerical_assessment}, with interchangeable conical nozzles of 1150 mm and 1950 mm diameter~\cite{scirocco_characterization}. The facility's research emphasizes large-scale capsule and vehicle testing under high-enthalpy conditions. Development and testing of Ultra High Temperature Ceramic nose caps coupling C/SiC materials with novel UHTC composites~\cite{scirocco_synergy} represents capsule-specific thermal protection research, examining aerothermal coupling effects where surface recession from ablation interacts with shock layer chemistry. France's ONERA operates the F4 facility at Le Fauga-Mauzac, a high-enthalpy hot shot wind tunnel. The facility provides flows at velocities from 5 to 6 km/s and reservoir pressures reaching several hundred atmospheres, with reservoir enthalpy capability up to 18 MJ/kg for air~\cite{f4_scramjet_tests, f4_validation}. Real gas effects are significant, with the free stream flow observed to be close to thermochemical equilibrium for pressure and translational temperature, while nitric oxide concentration measurements indicate near-nonequilibrium conditions~\cite{f4_flow_characterization}. Velocity measurements from Doppler shifts range between 2500 and 4500 m/s~\cite{f4_synthesis}. Research spans both fundamental transition studies and applied scramjet testing. Studies examined real gas influence on flare-induced boundary layer separation, addressing how high-temperature thermochemistry affects shock-boundary layer interactions~\cite{f4_validation}. The facility was extensively developed for scramjet combustion testing, with a 1.5-meter LAPCAT II small-scale flight experiment configuration tested at Mach 7.4 conditions, demonstrating positive aero-propulsive balance~\cite{f4_scramjet_tests}.





% Wind tunnels and their works: \texttt{VKI Longshot Hypersonic Tunnel},~\texttt{SCIROCCO}, \texttt{T6 Stalker Tunnel} Oxford,~\texttt{HIEST} JAXA, \texttt{JF-22} China,~\texttt{Mach 6 and 8 Quiet Wind Tunnel},~\texttt{Hypersonic Tunnel Facility (HTF),\ \gls{NASA}} 

% Furthermore, one of the biggest issues of experimental data is the acquisition of vibrational temperature in wind tunnels, this process 



% One of the greatest challenges in hypersonic flight research is the design of the \gls{TPS} and control surfaces, with most boundary layer transition research focussing on these two critical areas due to their exposure to the most extreme aerothermal environments. The material selection for \gls{TPS} is particularly demanding, as the extreme aerothermal environment fundamentally alters material behaviour and performance at the atomic and microstructural levels. \cite{PetersEtAl2024} Common metals and alloys such as aluminium and nickel-base superalloys are limited to moderate thermal loads below 800°C, whilst refractory metals can operate between 800-1200°C in oxidising atmospheres, and refractory ceramics can withstand temperatures exceeding 1700°C, though they lack monolithic thermal shock resistance. The heating intensity scales dramatically with velocity—a vehicle at Mach 5 experiences over 100 times the heating of one at Mach 1 at the same altitude, with a further threefold speed increase to Mach 15 amplifying heating by an additional factor of nearly 30. \cite{PetersEtAl2024} This extreme heating triggers multiple degradation mechanisms: at temperatures above 3000°C, molecular oxygen and nitrogen dissociate into highly reactive free radicals that dramatically accelerate oxidation rates, whilst carbonaceous composites begin oxidising at approximately 370°C with catastrophic degradation beyond 500°C. These coupled thermal, chemical, and mechanical effects impose stringent constraints on material selection and necessitate sophisticated protective coating systems to enable sustained hypersonic flight.

% The stringent material temperature limitations directly constrain the thermal boundary conditions experienced during hypersonic flight, with wall temperatures ranging from moderate (< 800°C for aluminium alloys) to extreme (> 1700°C for refractory ceramics) depending on material selection and vehicle design. These wall temperature variations have profound implications for boundary layer stability and transition behaviour, as thermal conditions fundamentally alter the instability mechanisms that govern laminar-to-turbulent transition. Accurate prediction of transition under realistic thermal loading conditions is therefore essential for vehicle design, yet experimental investigation of these effects presents significant challenges in ground-based testing facilities.



% Significant advances have been made in the last years in the study of high speed flows. Most detailed studies rely on CFD analyses, since carrying out physical experiments in the working conditions of interest is generally a costly or even infeasible task \citep{Bertin2006}. To get insights into the hypersonic regime, hyperveliocioty windtunnels have been designed to generate high speed flows in their working section. Unfortunately, the pwoer required to obtain the prescribed flow conditions is directly proporrtilnal to the flwo density, velocity, stagnationn enthalpy and cross-sectiopnal aread. Reproducing such severe conditions may cause overheating and result in structural damage; to avoid such problems, the test duration is often short. Moreover, instrumentation working efficiently at such extreme temperatures is also difficult to obtain. For all these reasons, hte cost of carrying the experiments are relatively high. However, there exist high ethalpy ground based facilities and in the flows we mention some of them. The NASA Hypersonic Tunnel Facility (HTF) originally designed to test nuclear thermal rockets, is capable of simulating up to Mach 7 at high enthalpy conditions. To simulate the effects of hypersonic conditions on TPS of re entry vehicles, the SCIROCCO and GHIBLI facilities of Plasma Wind Tunnel Complex at CIRA are able to reach Mach 12 and extremely high heat fluxes; the test chamber of SCIROCCO, in which the experimental measurements of pressure and temperature are carried out, is shown. In France, the PHEDRA wind tunnel simulates plasma gases of nitrogen air and methane and the ONERA F4 wint tuneell provides air flows at high enthlapy and pressure. The Plasmatron at Von Karman Institute is also a high enthalpy facility, which reaches temperatures up to 10000K; hot gas from the test chamber exits through a 700 kW heat exchanger and a 1050 kW cooling  system used. To overcome the aforementioned problematic, a number of hypersonic wind tunnels work at low enthalpy cryogenic conditions. The Purdue Mach 6 quiet tunnel, built by the end of 90s, is designed to achieve laminar nozzle-wall boundary layers to study natural transition to turbulence in cryogenic conditions. The Sandia Hypersonic wind tunnel accelerates the flow up to Mach 5 for an air mixture, and Mach 14 for a nitrogen mixture and the stagnation temperature is of the order of  1400K; this facility reproduces reentry vehicles or rockets flight conditions. A Mach number equal to 14 is reached by the AEDC Hypervelocity Wind Tunnel 9.



% One of the biggest challenges of hypersonic flight is the design of the \gls{TPS} and it's material selection. \cite{PetersEtAl2024} gives a comprehensive review of the material science ...  

% The practical challenges of hypersonic flight have been demonstrated through decades of experimental programs. \cite{national_air_space_museum_2012} Since the first hypersonic flight achieved by the Bumper-WAC vehicle in 1949 at 5,150 mph, numerous vehicles including the X-15, Space Shuttle, and more recently the X-43 and X-51 have achieved hypersonic speeds, though most shared a common limitation: after initial boost, they essentially glided through the atmosphere with only gravity as the propelling force. \cite{national_air_space_museum_2012, intechopen_hypersonic_2019} The X-43 achieved approximately 10 seconds of sustained scramjet thrust at Mach 9.6, while the X-51 operated for about 200 seconds at Mach 5, representing critical proof-of-concept demonstrations. \cite{national_air_space_museum_2012} More recently, successful flight tests of DARPA's Hypersonic Air-breathing Weapon Concept in 2021-2022 and the reusable Talon A2 vehicle exceeding Mach 5 in 2024-2025 mark the transition from fundamental research to engineering validation. \cite{wikipedia_scramjet_2025, spacenews_stratolaunch_2025} These experimental programs have consistently revealed that thermal protection and materials limitations represent the primary constraints on vehicle performance and operational envelope.





% In the 1980s, the U.S. launched the National Aerospace Plane (NASP) program to develop a hypersonic, reusable single-stage-to-orbit vehicle, with President Reagan envisioning flights from New York to Tokyo in two hours. \cite{japcc_hypersonic_vehicles_2022} The program encountered insurmountable technical challenges in scramjet propulsion, thermal protection, and materials that could withstand sustained hypersonic flight conditions. \cite{japcc_hypersonic_vehicles_2022} NASP was ultimately cancelled in 1992, though the decade of research provided foundational knowledge that informs current hypersonic vehicle development. \cite{japcc_hypersonic_vehicles_2022}




% \cite{PetersEtAl2024} In the last decade, there has been a resurgence in hypersonic vehicle development driven by the desire to increase flight performance and reusability. The maturation of scramjet propulsion systems, advances in ultra-high temperature ceramics, and expansion of ground-based testing infrastructure have fundamentally altered the feasibility of sustained atmospheric flight at speeds exceeding Mach 5. However, this progress has simultaneously revealed the profound complexity of hypersonic flow physics—particularly boundary layer transition, thermo-chemical nonequilibrium effects, and material oxidation at extreme temperatures. This section examines the key developments enabling contemporary low-altitude hypersonic vehicle programs, the fundamental material limitations constraining operational envelopes, and the critical role of wind tunnel and computational fluid dynamics capabilities in validating predictive models essential for vehicle design.






% Advances in hypersonic aerothermodynamics have been driven by renewed demand for sustained high-speed flight not only for reentry capsules but also for manoeuvrable \gls{HGV} and advanced air-breathing propulsion concepts. Historically, early vehicle designs relied on robust, heavily oversized \gls{TPS} and conservative margins due to the lack of reliable prediction tools for transition and thermal loads. Empirical design rules, developed from limited flight data and wind tunnel tests, often led to over-designed heat shields and higher structural weights to account for uncertainties.

% The fundamental understanding of hypersonic airflow physics and the establishment of comprehensive test programs represent perhaps the most critical knowledge gaps in contemporary aerospace research. The importance of this research domain extends far beyond academic curiosity—it directly determines whether hypersonic vehicles can transition from experimental demonstrators to operational systems, and whether theoretical predictions can be validated against physical reality.

% The fundamental understanding of hypersonic flow physics and the establishment of comprehensive literature perhaps the most critical knowledge gaps in contemporary aerospace research. 


% Advances in computational capabilities, experimental facilities, and materials science have steadily shifted this paradigm. Modern vehicles are now being designed for longer atmospheric flight times, tighter trajectory envelopes, and higher performance demands, all of which place greater emphasis on accurate thermal prediction. Material science developments — such as ultra-high-temperature ceramics, novel carbon–carbon composites, and oxidation-resistant coatings — have pushed operational temperature limits well beyond what was possible during earlier programs.

% However, even with improved materials, the maximum surface temperatures achievable still fall short of the adiabatic wall temperatures that can occur on leading edges and stagnation points in sustained hypersonic flight. As a result, the use of cold wall conditions remains an integral design strategy to protect underlying structures and systems. This means that reliable modelling tools are needed to quantify not only the overall thermal loads but also the critical locations where transition from laminar to turbulent flow might occur earlier than expected due to the interaction of surface temperature and boundary layer instabilities.

% Recent progress has focused on integrating advanced material capabilities with predictive design methods. Designers now face the challenge of leveraging these high-temperature materials efficiently while minimising mass and ensuring structural safety. Improved modelling tools that capture the multi-physics nature of hypersonic flows — including shock–boundary layer interactions, local heat fluxes, and realistic material behaviour under extreme conditions — are essential for closing the gap between material performance limits and real-world operational environments.

% At the same time, the shift towards reusable vehicles and more frequent hypersonic missions has raised the stakes for design accuracy. Conservative margins are no longer sufficient when even small discrepancies in predicted heating rates can lead to costly overdesign or, worse, catastrophic failure. Better validated models allow engineers to make informed trade-offs between material selection, active cooling strategies, and flight trajectory design, ultimately enabling the next generation of hypersonic systems to operate safely and efficiently within ever tighter constraints.








% \gls{DNS} of flow close to such extreme hypersonic regimes with thermochemical framework has recently appeared in the literature, with the first studies performed by \cite{marxen2013} on laminar boundary layers with instabilities suitable for hypersonic cruise conditions, whilst an extension study \citep{marxeniaccarinomagin2014} focused on non linear instability in the presence of finite rate chemistry without considering thermal non-equilibrium, concluding the dominant influence of chemical properties on primary disturbance amplitude, neglecting secondary instabilities during transition. Simulations of chemical non-equilibrium have received far greater attention as opposed to thermal non-equilibrium, with the \texttt{Mutation++} library \cite{Scoggins2020}, providing an extensive state of the art computation ability of physiochemical properties and has been widely used in the chemical studies concerning hypersonic flows. \cite{DiRenzo2020} introduced their \texttt{HTR} finite-difference solver for hypersonic aerothermodynamics for exascale computing, \cite{direnzourzay2021} used the routine to produce spatial evolution of such boundary layers to fully turbulent state. \cite{Passiatore2021} studied the complete spatially evolving turbulence transition process in thermochemical non-equilibrium for a quasi-adiabatic non-catalytic wall, where a second-mode transition was triggered with a strong amplitude and observed chemical activity to have little effect on the sub-harmonic transition. These studies were further extended to understand transition process in flat-plate boundary layer configurations with non-zero pressure gradients, i.e. oblique shock interactions by \citep{Volpiani2021} using another finite-difference code \texttt{CHARLIE}, and \cite{passiatore2023}, which is essentially an thermochemical extension of the shock wave interaction work of \cite{sandham2014} in hypersonic cold flow. The studies have concluded that mean velocities and turbulent intensity profiles are similar to those of cold hypersonic flows, in fact, the presence of varying species mass fractions caused a reduction in the turbulent fluctuations. At this point,  a complete understanding of chemical and thermal influences on modal stability for high Reynolds number, high Mach number flows is available in the literature. 



% To date, various CFD solvers have been developed by some universities and scientific
% research institutions for nonequilibrium flows in the near-continuum regime. These
% include some proprietary solvers, such as NASA Langley code LAURA [58], NASA
% Ames code DPLR

% \gls{NASA} Langley code \texttt{LAURA} \citep{Cheatwood1996}
% NASA AMES code \texttt{DLPR} \citep{Candler1994}
% \texttt{LeMANS} solver Michigan \citep{Scalabrin2005}
% \texttt{US3D} Minnesota \citep{Nompelis2004}
% \texttt{DLR-TAU} \citep{Mack2002}

% \texttt{Eilmear} \citep{Gollan2013}
% \texttt{SU2-NEMO} \citep{MaierEtAl2021}, 
% \texttt{HTR solver} \citep{DiRenzoFuUrzay2020}
% \texttt{hy2Foam} \citep{Casseau2016a,Casseau2016b}

% \texttt{CEA}~\citep{McBrideGordon1996}
% \texttt{Cantera}~\citep{GoodwinMoffatSpeth2015}
% \texttt{KAPPA}~\citep{CampoliEtAl2019}
% \texttt{Mutation++}~\citep{Scoggins2020}

% The route to turbulence in hypersonic boundary layers is not unique; rather, breakdown occurs through different physical mechanisms that depend on the dominant instability modes, disturbance spectrum, and surface conditions.


% \newpage

% \section{Scope of the present work}

% This thesis investigates boundary layer transition mechanisms in hypersonic flows through high-fidelity Direct Numerical Simulation. The research focuses on second-mode instability development and breakdown on canonical geometries at conditions representative of low-altitude, high-speed flight where ground facilities face significant limitations. A compressible Navier-Stokes solver with finite-rate chemistry and thermal nonequilibrium capabilities has been developed specifically for this work, enabling detailed examination of instability growth, nonlinear interactions, and transition physics under realistic thermochemical conditions. The primary objectives are to characterise second-mode behaviour at elevated enthalpies, quantify the influence of thermochemical nonequilibrium on transition, and provide high-fidelity data for validation of reduced-order transition prediction methods. Simulations are performed on sharp and blunted cones at Mach numbers between 6 and 10, covering conditions from calorically perfect gas to significant real-gas effects. This work contributes fundamental understanding of hypersonic transition physics whilst demonstrating computational capabilities for parameter studies beyond experimental reach.



\section{Scope of the Present Work}\label{section:introduction/scope}

The preceding sections have outlined the fundamental physics of hypersonic boundary layer transition, documented capabilities and limitations of experimental facilities, and reviewed computational approaches for simulating high-enthalpy flows. Whilst significant progress has been made in understanding second-mode instability under idealised conditions, transition behaviour under combined effects of high enthalpy, finite-rate chemistry, and realistic flow conditions remains incompletely characterised. Existing experimental data are constrained by facility limitations, and many numerical studies have employed simplified thermochemical models or addressed individual effects in isolation. This thesis aims to address specific gaps in understanding hypersonic transition physics through high-fidelity Direct Numerical Simulation, with particular emphasis on thermochemical nonequilibrium effects and instability behaviour at conditions representative of low-altitude, high-speed flight. The following paragraphs outline the specific objectives, geometric configurations, flow conditions, and computational methodology employed in this work, along with justification for key modelling decisions.

% Accurately predicting the transition from laminar to turbulent flow in hypersonic boundary layers remains one of the most challenging and consequential problems in high-speed aerothermodynamics. While linear stability theory (LST) has greatly improved our understanding of the initial growth of instabilities — particularly the second-mode (Mack mode) acoustic waves that dominate at higher Mach numbers — the detailed pathways through which these disturbances interact and break down into turbulence are still not fully understood.

Transition to turbulence in hypersonic boundary layers has been studied in the past, both in \textit{`cold flow'} and high-enthalpy regime at suborbital enthalpies, despite the challenges imposed by the later environment. The studies concern stages of transition through perturbations in the laminar base flow through linear and non-linear stability analyses. \cite{Cerminara2019, cerminara_sandham_2020} studied the receptivity mechanisms at the leading edge of a blunt body in cold hypersonic flow without considering finite-rate chemical models. Studies of transition and control for chemically perfect gases, and flows in chemical equilibrium were first performed by  \citep{Malik1989, Malik1990b, Malik1990a, Malik1990c}. This was extended further by \cite{Stuckert1994} and \cite{Hudson1997}, who introduced chemical and vibrational non-equilibrium extensions with Park's two-temperature model. The studies conclusively show that high enthalpy effects bring non-negligible contribution to boundary layer stability; the endothermic nature of the reactions remove energy from the thermal boundary layer, leaving the flow more subjected to second-mode (Mack modes) disturbances, almost mimicking wall-cooling effects. \cite{Malik2003} and \cite{Johnson2005}, conducted parabolised stability equations (PSE) to the subject, with the former study concluding that some dissociation reactions actually reduce the growth of perturbations caused by the chemical reactions on the base flow. \cite{Johnson2005} has also performed simulations using multi-species models to better capture the stability during transition. \cite{Franko2010} concluded that the chemical and viscosity model used in the analysis has a substantial effect on the stability characteristic. Consequently, \cite{Klentzman2013} showed that viscosity stabilises the base flow while investigating viscous effects on the receptivity of boundary layers and \cite{miromiro2018} extended the work on an adiabatic flat plate, observing the flow to most destabilise when mass diffusion is hundreds of times faster than viscous diffusion and further investigated the effects of choice of thermochemical and thermophysical framework on the transition prediction \citep{MiroMiro2019}. Furthermore, work of \cite{wartmann2018} show that thermochemical non-equilibrium approach for the mean flow and the stability code are almost identical to the approach of considering just the base flow in thermochemical non-equilibrium. Most of the work in the literature focused on the modal growth of the transition process predominantly for sub-orbital enthalpies. The stability of high-speed flows at low altitudes, where the Reynolds numbers are significantly higher and a comprehensive study of the evolution of the boundary layer to complete turbulence transition in such flows has been largely unexplored areas.  

% Modal growth comparisons
% This uncertainty is amplified in conditions relevant to modern hypersonic vehicles, where high temperatures and complex thermodynamic effects significantly alter the evolution of boundary layer instabilities. The consequences of transition prediction errors are substantial: premature transition can lead to localised heat loads well beyond what thermal protection systems are designed to handle, jeopardising structural integrity and mission safety. Advances in material science have pushed the operational temperature limits higher than ever before, but the need for surface cooling remains, which can further destabilise certain modes and complicate transition prediction.

Direct numerical simulations (DNS) of flow close to such extreme hypersonic regimes with thermochemical framework has recently appeared in the literature, with the first studies performed by \cite{marxen2013} on laminar boundary layers with instabilities suitable for hypersonic cruise conditions, whilst an extension study \citep{marxen_iaccarino_magin_2014} focused on non linear instability in the presence of finite rate chemistry without considering thermal non-equilibrium, concluding the dominant influence of chemical properties on primary disturbance amplitude, neglecting secondary instabilities during transition. Simulations of chemical non-equilibrium have received far greater attention as opposed to thermal non-equilibrium, with the Mutation++ library \cite{Scoggins2020}, providing an extensive state of the art computation ability of physiochemical properties and has been widely used in the chemical studies concerning hypersonic flows. \cite{DiRenzo2020} introduced their HTR finite-difference solver for hypersonic aerothermodynamics for exascale computing, \cite{direnzo_urzay_2021} used the routine to produce spatial evolution of such boundary layers to fully turbulent state. \cite{Passiatore2021} studied the complete spatially evolving turbulence transition process in thermochemical non-equilibrium for a quasi-adiabatic non-catalytic wall, where a second-mode transition was triggered with a strong amplitude and observed chemical activity to have little effect on the sub-harmonic transition. These studies were further extended to understand transition process in flat-plate boundary layer configurations with non-zero pressure gradients, i.e. oblique shock interactions by \citep{Volpiani_2021} using another finite-difference code \textit{CHARLIE}, and \cite{passiatore2023}, which is essentially an thermochemical extension of the shock wave interaction work of \cite{sandham2014} in hypersonic cold flow. The studies have concluded that mean velocities and turbulent intensity profiles are similar to those of cold hypersonic flows, in fact, the presence of varying species mass fractions caused a reduction in the turbulent fluctuations. At this point,  a complete understanding of chemical and thermal influences on modal stability for high Reynolds number, high Mach number flows is available in the literature. 


writing a DNS solver to answer the questions

Despite this progress, key questions remain about how second-mode instabilities evolve in different thermodynamic environments and how they interact with other disturbance modes to form turbulent spots. A purely linear perspective cannot capture these interactions; therefore, a combined approach that integrates LST with DNS offers a more complete picture. Using LST to guide the identification of dominant modes and interpret their initial growth provides critical context for understanding how these instabilities develop into fully turbulent flow when all relevant physics are included.

% How everything sums up together? 
% This project aims to address these knowledge gaps by applying state-of-the-art DNS, complemented by LST analysis, to investigate the breakdown mechanisms of Mack modes in high-speed boundary layers under a range of thermal conditions. By clarifying how these instabilities grow, interact, and trigger transition, this work will contribute to the development of more reliable transition prediction methods. Such advances are essential for designing robust, efficient thermal protection systems and reducing uncertainty in the performance of next-generation hypersonic vehicles.


% \newpage

\subsection{Aims and objectives}\label{subsection:aimsandobjectives}

This project aims to address these knowledge gaps by applying state-of-the-art DNS, complemented by LST analysis, to investigate the breakdown mechanisms of Mack modes in high-speed boundary layers under a range of thermal conditions. By clarifying how these instabilities grow, interact, and trigger transition, this work will contribute to the development of more reliable transition prediction methods. Such advances are essential for designing robust, efficient thermal protection systems and reducing uncertainty in the performance of next-generation hypersonic vehicles.

Following the motivation outlined above, this project aims to develop a rigorous mathematical and computational framework for analysing high-enthalpy hypersonic boundary layers under thermochemical non-equilibrium conditions, alongside a physical understanding of transition to turbulence. Ionisation effects are not considered in this study. The research first seeks to extend self-similar boundary-layer theory, which has received relatively little attention in the low-altitude, high-enthalpy regime where multiple species and nonequilibrium effects play a significant role. While most \gls{LST} studies in the literature assume a frozen freestream base flow~\citep{miromiroThesis}, the post-shock region near the stagnation point is typically closer to local thermochemical equilibrium, which remains poorly understood. Furthermore, the project investigates the transition mechanisms in high-Mach-number, high-Reynolds-number conditions under varying wall boundary conditions.

To address these challenges, part of the work focuses on extending the in-house finite difference compressible solver OpenSBLI to incorporate detailed thermal and chemical nonequilibrium models with appropriate shock-capturing techniques. The project also aims to clarify how extreme nonlinear interactions drive transition to turbulence under these conditions. Taking into account the current state of the art, the following objectives are proposed to address this critical gap in our understanding of transition on a wedge and the leading edge of a blunt body in thermochemical nonequilibrium.

% [topsep=0pt,itemsep=-1ex,partopsep=1.5ex,parsep=3.5ex]

\begin{enumerate}
    \item To extend the analytical set of Fay-Riddell equations to include multiple species,  thermochemical non-equilibrium, and Park's two-temperature model.

    \item Extend the in-house compressible solver OpenSBLI to include thermochemical framework with chemical and vibrational non-equilibrium, including validation for flow around a cylinder to understand the stagnation point behaviour. 
    

    \item To investigate the effects of LST and its application in e\textsuperscript{N} method to accurately predict the behaviour of instability modes and their sensitivities to the high-enthalpy base flow inaccuracies.

    
    \item Understand the mechanics of temporally and spatially developing turbulent flow region, using DNS of spreading turbulent spots and wedge of developing turbulent flow. 

\end{enumerate}


% \newpage

\section{Thesis outline}
The structure of this report is as follows: \Cref{chapter:governing_equations} presents the necessary theoretical framework and governing equations used to describe the fundamental flow physics, including the thermochemical processes relevant to high-enthalpy boundary layers. \Cref{chapter:numerical_methods} introduces the \gls{DNS} solver and details the numerical methods employed, along with representative validation cases that demonstrate the solver’s ability to resolve the complex features of transitional hypersonic flows. \Cref{chapter:disturbance_growth} examines the propagation and amplification of small disturbances leading to second-mode breakdown in thermally excited boundary layers, drawing systematic comparisons between \gls{LST} and \gls{DNS} results to highlight the role of vibrational nonequilibrium and real gas effects. Building on this, \Cref{chapter:breakdown} focuses on isolating and analysing Mack modes to clarify the mechanisms by which these dominant instability modes undergo complete breakdown and trigger transition under various thermally excited conditions. Finally, \Cref{Chapter: Conclusions} summarises the main findings, discusses their broader implications for hypersonic transition prediction and \gls{TPS} design, and provides recommendations for future research directions.